%% Delete this \nocite command invocation to make the references section only list out the bibitems that are actually cited.


\nocite{*}

\section{The problems}

\begin{definition} \label{definition:superelliptic_curve_over_a_field}
    Let $k$ be a \CrefAndHyperrefIfExist{definition:field}{field}. 
    \TODO{projective closure, plane curve, affine equation}
    \TODO{smooth completion of an affine curve}
    A \hldef{superelliptic curve over $k$} is the smooth completion of the plane curve with affine equation of the form 
    $$y^d = f(x)$$
    where $f(x) \in k[x]$ and $d \geq 2$ is some integer.
\end{definition}

\begin{definition} \label{definition:mu_d_twist_of_a_superelliptic_curve_by_a_scalar}
    Let $k$ be a \CrefAndHyperrefIfExist{definition:field}{field}. Let $C$ be a \CrefAndHyperrefIfExist{definition:superelliptic_curve_over_a_field}{superelliptic curve} over $k$ given as the smooth completion of the plane curve with affine equation 
    $$y^d = f(x)$$
    where $f(x) \in k[x]$ and $d \geq 2$ is some integer.
    Given an element $u \in k^\times$, the \hldef{$\mu_d$-twist of $C$} is the superelliptic curve given as the projective closure of the plane curve with affine equation
    $$y^d = \frac{1}{u} f(x).$$
\end{definition}

Here is the main problem:

\begin{context} \label{context:main_context_for_average_dimensions_of_selmer_groups_of_twists_of_superelliptic_curve_over_global_function_field}
    Let $K = \Fq(t)$ be a \CrefAndHyperrefIfExist{definition:global_function_field}{global function field} of \CrefAndHyperrefIfExist{definition:characteristic_a_prime_for_a_scheme}{characteristic} coprime to $6$. Let $\ell \geq 5$ be a prime number such that $\gcd(\ell, q) = 1$. Write $F_{n,\ell}$ for the set of $\ell$th power free polynomials over $\Fq$ of degree $n$. Let $v$ be a prime number coprime to $6 q \ell$. 
\end{context}

\begin{context} \label{context:mu_ell_twists_of_superelliptic_curve_over_global_function_field}
    Assume \Cref{context:main_context_for_average_dimensions_of_selmer_groups_of_twists_of_superelliptic_curve_over_global_function_field}. Fix a \CrefAndHyperrefIfExist{definition:superelliptic_curve_over_a_field}{superelliptic curve} $C/K$ whose affine equation is given by 
    $$C : y^\ell = f(x)$$
    for some non-trivial polynomial $f(x) \in K[x]$. 
    
    Given an $\ell$-th power free polynomial $g \in \Fq[t]$ write \hl{$C_g$} for the \CrefAndHyperrefIfExist{definition:mu_d_twist_of_a_superelliptic_curve_by_a_scalar}{$\mu_\ell$-twist} of $C$ by $g$. In other words $C_g$ is the smooth completion of the plane curve with affine equation of the form 
    $$C_g: g y^\ell = f(x)$$
\end{context}

\begin{problem} \label{problem:main_average_dimensions_of_selmer_groups_of_twists_of_superelliptic_curve_over_global_function_field}
    % Let $K = \Fq(t)$ be a \CrefAndHyperrefIfExist{definition:global_function_field}{global function field} of \CrefAndHyperrefIfExist{definition:characteristic_a_prime_for_a_scheme}{characteristic} coprime to $6$. Let $\ell \geq 5$ be a prime number such that $\gcd(\ell, q) = 1$. Write $F_{n,\ell}$ for the set of $\ell$th power free polynomials over $\Fq$ of degree $n$. Let $v$ be a prime number coprime to $6 q \ell$. 
    Assume \Cref{context:main_context_for_average_dimensions_of_selmer_groups_of_twists_of_superelliptic_curve_over_global_function_field} and \Cref{context:mu_ell_twists_of_superelliptic_curve_over_global_function_field}

    For each integer $j \geq 0$, what is the asymptotic value of the quantity
    $$\frac{\# \{g \in F_{n,\ell}:  \dim_{\bbF_v} \Sel_v (\Jac(C_g) / K ) = j \} }{\#F_{n,\ell}}$$
    (\Cref{definition:selmer_group_for_an_isogeny_of_commutative_group_schemes_over_a_global_field}) (\Cref{definition:jacobian_of_a_curve_over_a_field})
    as $n \to \infty$?
\end{problem}



Here are subproblems

\begin{problem}
\end{problem}

\begin{problem}
    Assume \Cref{context:main_context_for_average_dimensions_of_selmer_groups_of_twists_of_superelliptic_curve_over_global_function_field}.

    Write $\mathcal{J}_g$ for the N\'eron model of $\Jac(C_g)$ over 
\end{problem}


\section{Hall's big monodromy approach}

\input{../_concepts/theorem_langs_theorem_for_connected_algebraic_groups_over_finite_fields.tex}
\begin{definition} \label{definition:kummer_sheaf_associated_to_a_character_on_points_of_a_connected_commutative_algebraic_group_over_a_finite_field}
    Let $G/\Fq$ be a conected commutative algebraic group. Let $\ell \neq \Char \Fq$ be a prime number. Let $n \geq 1$.
    By \CrefAndHyperrefIfExist{theorem:langs_theorem_for_connected_algebraic_groups_over_finite_fields}{Lang's theorem}, given a connected commutative algebraic group $G/\Fq$, we have we have a short exact sequence
    $$1 \to G(\Fqn) \to G_{\Fqn} = G \times_{\Fq} \Fqn \xrightarrow{L} G_{\Fqn} \to 1$$
    where $L$ is the \CrefAndHyperrefIfExist{definition:lang_map_of_an_endomorphism_of_an_algebraic_group}{Lang isogeny}. Since this identifies $G_{\Fqn}$ with a \CrefAndHyperrefIfExist{definition:finite_etale_cover_of_a_scheme}{finite \'etale} $G(\Fqn)$-cover of $G_{\Fqn}$, we have a \CrefAndHyperrefIfExist{theorem:finite_etale_covers_of_a_connected_scheme_correspond_to_finite_sets_with_an_action_of_the_etale_fundamental_group}{corresponding} surjective group homomorphism $\pioneet(G_{\Fqn}, e) \to G(\Fqn)$. 
    
    \begin{enumerate}
        \item Let $\chi: G(\Fqn) \to \bbF_\ell^\times$ be a character. There is a \CrefAndHyperrefIfExist{definition:lisse_sheaf_on_a_scheme}{lisse sheaf} \hl{$\calL_\chi$} of rank $1$ with coefficients in $\bbF_\ell$ on $G_{\Fqn} = G \times_{\Fq} \Fqn$ \CrefAndHyperrefIfExist{theorem:lisse_ell_adic_sheaves_on_connected_noetherian_schemes_are_equivalent_to_representations_of_etale_fundamental_group_over_adic_ring}{corresponding} to the composed group homomorphism
        $$\pioneet(G_{\Fqn}, e) \to G(\Fqn) \xrightarrow{\chi} \bbF_\ell^\times.$$
        We may call this lisse sheaf the \hldef{Kummer sheaf associated to $\chi$}.

        Similarly, Kummer sheaves are associated to characters $\chi: G(\Fqn) \to \bbZ_\ell^\times$, $\chi: G(\Fqn) \to \bbQ_\ell^\times$, $\chi: G(\Fqn) \to \overline{\bbQ}_\ell^\times$, etc.; each type of Kummer sheaf is a lisse sheaf on $G_{\Fqn}$ with appropriate coefficients.
    \end{enumerate}

\end{definition}

\begin{definition} \label{definition:decomposition_and_inertia_subgroups_of_etale_fundamental_group_of_dense_open_of_integral_normal_scheme_at_point_in_complement}
    Let $X$ be an \CrefAndHyperrefIfExist{definition:integral_scheme}{integral} \CrefAndHyperrefIfExist{definition:normal_and_factorial_scheme}{normal scheme} with \CrefAndHyperrefIfExist{definition:generic_point_of_a_topological_space}{generic point} $\eta$ and \CrefAndHyperrefIfExist{definition:function_field_over_a_field}{function field} $K$. Let $\overline{\eta}$ be a geometric point of $X$ whose image is $\eta$. Let $U \subseteq X$ be a non-empty \CrefAndHyperrefIfExist{definition:open_and_closed_subscheme_of_a_scheme}{open subscheme}. Let $x \in X \setminus U$ be a point. Let $v$ be a \CrefAndHyperrefIfExist{definition:valuation_on_a_field}{valuation} of $K$ \CrefAndHyperrefIfExist{definition:centered_at_for_a_prime_of_an_integral_domain_point_of_an_integral_scheme}{centered at} $x$\CrefIfExists{proposition:integral_scheme_has_a_valuation_centered_at_every_point_ideal}. 
    

    Recall that there is a natural surjection
    $$\operatorname{Gal}(K^{\text{sep}} / K) \twoheadrightarrow \pioneet(U, \overline{\eta})$$
    (\Cref{proposition:mor_from_et_fund_grp_of_gen_pt_of_norm_int_sch_to_et_fund_grp_of_scheme_is_quo_map_of_gal_grp_of_sep_clos_of_func_field_to_gal_grp_of_max_unr_ext})

    \begin{enumerate}
        \item The \hldef{decomposition subgroup} \hl{$D_v$} of $\pioneet(U, \overline{\eta})$ is the image of the \CrefAndHyperref{definition:decomposition_group_of_a_general_galois_extension_of_valuation_fields}{decomposition subgroup} $D_{K,v}$ of $\operatorname{Gal}(K^{\text{sep}} / K)$ in $\pioneet(U, \overline{\eta})$.

        \item The \hldef{inertia subgroup} \hl{$I_v$} of $\pioneet(U, \overline{\eta})$ is the image of the \CrefAndHyperref{definition:inertia_group_of_a_general_galois_extension_of_valuation_fields}{inertia subgroup} $I_{K,v}$ of $\operatorname{Gal}(K^{\text{sep}} / K)$ in $\pioneet(U, \overline{\eta})$.
    \end{enumerate}
    Given two choices $v_1$ and $v_2$ of valuations centered at $x$, the groups $D_{v_1}$ and $D_{v_2}$ and the groups $I_{v_1}$ and $I_{v_2}$ are conjugate to each other in $\pioneet(U, \overline{\eta})$. As such, one might write \hl{$D_x$} and \hl{$I_x$} for the conjugacy classes. 
\end{definition}
\begin{definition} \label{definition:ramification_of_an_etale_sheaf_on_a_connected_normal_scheme_at_a_geometric_point}
    Let $X$ be an \CrefAndHyperrefIfExist{definition:connected_object_of_a_galois_category}{connected} \CrefAndHyperrefIfExist{definition:normal_and_factorial_scheme}{normal scheme} with \CrefAndHyperrefIfExist{definition:generic_point_of_a_topological_space}{generic point} $\eta$ and \CrefAndHyperrefIfExist{definition:function_field_over_a_field}{function field} $K$. Let $\overline{\eta}$ be a geometric point of $X$ whose image is $\eta$. Let $U \subseteq X$ be a non-empty \CrefAndHyperrefIfExist{definition:open_and_closed_subscheme_of_a_scheme}{open subscheme}. Let $x \in X \setminus U$ be a point. 

    Let $\calF$ be one of the following:
    \begin{enumerate}
        \item A \CrefAndHyperrefIfExist{definition:cardinal_number_and_cardinality_of_a_set}{finite} \CrefAndHyperrefIfExist{definition:locally_constant_sheaf}{locally constant sheaf} of \CrefAndHyperrefIfExist{definition:zermelo_fraenkel_set_theory}{sets} on $U_{\et}$ 
        \item A finite locally constant sheaf of \CrefAndHyperrefIfExist{definition:cardinal_number_and_cardinality_of_a_set}{finite cardinality} \CrefAndHyperrefIfExist{definition:group}{abelian groups} on $U_{\et}$ 
        \item A locally constant sheaf of \CrefAndHyperrefIfExist{definition:finitely_generated_modules_over_rings}{finitely generated} \CrefAndHyperrefIfExist{definition:module_over_a_sheaf_of_rings_on_a_site}{modules} over a finite \CrefAndHyperrefIfExist{definition:ring}{ring} on $U_{\et}$ 
        \item In the case that $U$ is \CrefAndHyperrefIfExist{definition:irreducible_topological_space}{irreducible} and \CrefAndHyperrefIfExist{definition:unibranch_geometircally_unibranch_scheme}{geometrically unibranch}, a \CrefAndHyperrefIfExist{definition:finite_type_morphism_of_schemes}{finite type}, locally constant sheaf of \CrefAndHyperrefIfExist{definition:module_over_a_sheaf_of_rings_on_a_site}{$\Lambda$-modules} on $U_{\et}$
        \item Let $(R, \lambda)$ be a \CrefAndHyperrefIfExist{definition:complete_ring_with_respect_to_an_ideal}{complete} \CrefAndHyperrefIfExist{definition:discrete_valuation_ring}{DVR} such that $R/\lambda^n$ is finite for all $n \geq 1$. Let $E$ be the \CrefAndHyperrefIfExist{definition:field_of_fractions_of_an_integral_domain}{fraction field} of $R$. 
        \begin{enumerate}
            \item A lisse object of $\Sh(U,R)$ (\Cref{definition:lambda_adic_sheaf_on_the_small_etale_site_of_a_scheme_for_an_ideal_of_a_commutative_ring})
            \item A lisse object of $\Sh(U,E)$ (\Cref{definition:category_of_etale_sheaves_with_rational_coefficients_in_the_fraction_field_of_a_complete_dvr})
            \item For an \CrefAndHyperrefIfExist{definition:algebraic_element_minimal_polynomial_of_an_algebraic_element_algebraic_field_extension_transcendental_element_field_extension}{algebraic extension} $K$ of $E$, a lisse object of $\Sh(U,K)$ (\Cref{definition:category_of_etale_sheaves_with_rational_coefficients_in_an_algebraic_extension_of_the_fraction_field_of_a_complete_dvr}).
        \end{enumerate}
    \end{enumerate}

    In each of the above cases for $\calF$, there is a corresponding action of $\pioneet(U, \overline{\eta})$ on the stalk $\calF_{\eta}$ 
    (\Cref{theorem:locally_constant_sheaves_on_X_etale_are_equivalent_to_representations_of_the_etale_fundamental_group})
     (\Cref{theorem:lisse_ell_adic_sheaves_on_connected_noetherian_schemes_are_equivalent_to_representations_of_etale_fundamental_group_over_adic_ring}):
    $$\pioneet(U, \overline{\eta}) \to \Aut(\calF_{\eta})$$
    where $\Aut(\calF_{\eta})$ is the automorphism group for the appropriate category that $\calF_{\eta}$ belongs to. Further recall that there is an \CrefAndHyperrefIfExist{definition:decomposition_and_inertia_subgroups_of_etale_fundamental_group_of_dense_open_of_integral_normal_scheme_at_point_in_complement}{inertia subgroup} $I_x$ of $\pioneet(U, \overline{\eta})$, well defined up to conjugation. 

    We say that $\calF$ is \hldef{unramified at $x$} if the action of $I_x$ on $\calF_{\overline{\eta}}$ is trivial, i.e. the composed homomorphism
    $$I_x \hookrightarrow \pioneet(U, \overline{\eta}) \to \Aut(\calF_{\eta})$$
    is trivial. Otherwise, we say that $\calF$ is \hldef{ramified at $x$}.
\end{definition}

Let $C/\bbF_q$ be a smooth, geometrically connected curve. Let $X/K(C)$ be a superelliptic curve with affine model $y^r = f(x)$ for a polynomial $f(x) \in K(C)[x]$. Given $g \in K(C)^\times$, let $X_g$ be the twist of $X$ with affine model $g \cdot y^r = f(x)$. Write $\calX_g \to C$ for the N\'eron model of $X_g$. For a prime number $v$ not dividing $\Char \Fq$, $\Pic^0_{\calX_g/C}[v]$ is a quasi-finite \'etale group scheme over $C$ \TODO{cite}. 

% Let $\Conf_{X/C}^n$ be the configuration space of $n$ unordered points of $C$. 

\begin{context} \label{context:quadratic_twist_of_a_family_of_elliptic_curves_by_a_universal_family_of_polynomials}
    Let \hl{$q$} be a prime power. Fix an elliptic curve \hl{$E_1/\Fq(t)$} with non-constant $j$-invariant and let \hl{$\mathcal{E}_1 \to \bbP^1_{\Fq}$} be its \CrefAndHyperrefIfExist{definition:neron_model_of_an_abelian_variety_over_the_function_field_of_an_integral_scheme.tex}{N\'eron model}. Let \hl{$\ell$} be a prime not dividing $q$. The kernel $\calE_{1}[\ell]$ is a \CrefAndHyperrefIfExist{definition:etale_morphism_of_schemes}{\'etale} group scheme over $C$\CrefIfExists{proposition:multiplication_by_ell_for_commutative_group_scheme_over_general_base_is_etale_if_characteristic_of_residue_fields_is_different_from_ell} and in particular is \CrefAndHyperrefIfExist{definition:quasifinite_morphism_of_schemes}{quasi-finite}\CrefIfExists{proposition:etale_morphisms_of_schemes_are_quasifinite}.

    For a positive integer $d$, and for an algebraically closed extension field $\bark$ of $\Fq$, identify a geometric point $(C_0,\ldots,C_{d-1}) (\Spec \Fq[c_0,\ldots,c_{d-1}])(\bark)$ for $(C_0,\ldots,C_{d-1}) \in \bark^d$ with a monic polynomial $t^d + C_{d-1} t^{d-1} + \cdots + C_1 t + C_0 \in \Fqbar[t]$. Fixing a polynomial $m \in \bbF_q[t]$, let \hl{$F_d \subset \Spec \Fq[c_0,\ldots,c_{d-1}]$} be the parameter space of square-free polynomials of degree $d$ relatively prime to $m$. Explicitly, 
    $$F_d = \Spec \Fq[c_0,\ldots,c_{d-1}] \setminus (\Delta \cup V(m))$$
    \TODO{resultant}
    where \hl{$\Delta$} is the \CrefAndHyperrefIfExist{definition:discriminant_of_a_polynomial_over_a_commutative_ring}{discriminant} locus in $\Spec \Fq[c_0,\ldots,c_{d-1}]$ and \hl{$V(m)$} is the locus of polynomials that are not coprime with $m$. For each geometric point of $F_d(\Fqbar)$ corresponding to a polynomial $f \in \Fqbar[t]$, let \hl{$E_f/K$} be the \CrefAndHyperrefIfExist{definition:quadratic_twist_of_an_elliptic_curve_over_a_field}{quadratic twist} of by $f$. Let \hl{$\calE_f \to \bbP_\Fq^1$} be its \CrefAndHyperrefIfExist{definition:neron_model_of_an_abelian_variety_over_the_function_field_of_an_integral_scheme.tex}{N\'eron model}. 

    Let \hl{$S \coloneq \bbP_\Fq^1 \times F_d$}; in accordance to the notations above, it has coordinates $(t,c)$. Let \hl{$P(t,c) = t^d+c_{d-1} t^{d-1} + \cdots + c_1 t + c_0 \in \Fq[t][c_0,\ldots,c_{d-1}]$} be the universal monic polynomial of degree $d$ over $F_d$ --- it is a function on $\bbA^1 \times F_d \subset S$. Let \hl{$\calE_{\text{univ}} = \calE_1 \times_{\bbP_\Fq^1} S \to S$}. Let \hl{$\widetilde{S} \subset \bbA_\Fq^1 \times S$}, with coordinates $(y,t,c)$, be the closed subscheme defined by 
    $$w^2 = P(t,c).$$
    The natural projection 
    $$\hlin{\sigma: \widetilde{S} \to S, \quad (w,t,c) \mapsto (t,c)}$$
    is a finite flat double cover. Let \hl{$U_{\text{ram}} \subset S$} be the branch locus of $\sigma$, explicitly given by $P(t,c) = 0$. Let \hl{$U = S \setminus U_{\text{ram}}$}. In particular, the restriction $\sigma|_{\widetilde{S}|_U} : \widetilde{S}|_U \to U$ is a finite \'etale Galois cover with Galois group $\Z/2\Z$.

    Let \hl{$\iota: \widetilde{S} \to \widetilde{S}$} be the involution $(w,t,c) \mapsto (-w,t,c)$. Writing $(Q, w, t, c)$ for a point of 
    $$\sigma^* \calE_{\text{univ}} = \calE_1 \times_S \widetilde{S} = (\calE_1 \times_{\bbP_\Fq^1} S) \times_S \widetilde{S} = \widetilde{E}_1\times_{\bbP_{\Fq}^1} \widetilde{S},$$
    there is an action by $\bbZ/2\bbZ$ on $\sigma^* \mathcal{E}_{\text{univ}}$ by the action
    $$(Q, w, t, c) \mapsto (-Q, \iota(w,t,c)) = (-Q, -w, t, c).$$
    This action is \'etale and free on the generic fiber. The quotient $\sigma^* \calE_{\text{univ}} / (\bbZ/2\bbZ)$ of $\sigma^* \calE_{\text{univ}}$ by this action, which exists by faithfully flat descent, is a family of generalized elliptic curves $S$. 
    \TODO{generalized elliptic curve}
    \TODO{understand how faithfully flat descent works}
\end{context}

\begin{proposition}  \label{proposition:fiber_of_quadratic_twist_of_family_of_neron_model_of_elliptic_curves_over_F_q_t_over_polynomials_at_a_polynomial_is_quadaratic_twist_of_elliptic_curve_by_polynomial}
    Assume \Cref{context:quadratic_twist_of_a_family_of_elliptic_curves_by_a_universal_family_of_polynomials}.

    \TODO{should I be expressing this as a fiber of a point or a fiber over $\bbP_{\Fq}^1 \times \text{point}$? Because depending on whether I'm going for a point or a closed locus, I would be taking about a quadratic twist of $E_1$ or $\calE_1$.}
    For every geometric point $\overline{f} \in F_d(\Fqbar)$ corresponding to a square-free polynomial $f(t) \in \Fqbar[t]$, the restriction of $\calG \coloneq (\sigma^* \calE_{\text{univ}})/(\bbZ/2\bbZ)$, which is a family of generalized elliptic curves over $S$, to the fiber $\bbP_{\Fqbar}^1 \times \{\overline{f}\} \subset S$ is isomorphic to the N\'eron model $\calE_f \to \bbP_{\Fqbar}^1$ of the quadratic twist $E_f$. 
    \TODO{recheck proof, evaluating why it's the smooth locus}
    % The fiber of $(\sigma^* \calE_{\text{univ}})/(\bbZ/2\bbZ) \to S$ over a 
    % point $(t_0, C) \in S(\Fq)$ is the quadratic twist of $\mathcal{E}_1$ by $P(t_0,C)$. 
\end{proposition}
\begin{proof}
    Since $2$ is invertible in $\Fq$, the N\'eron model $\calE_1 \to C$ admits a relative Weierstrass equation
    % Note that an arbitrary elliptic scheme does not automatically admit a global Weierstrass equation, but we still have local Weierstrass equations.
    $$y^2 = x^3 + Ax + B$$
    \TODO{probably a relative Riemann-Roch theorem yields this}
    locally over $C$ where $A$ and $B$ are sections of $\underline{\omega}_{\calE_1/C}^{\oplus 4}$ and $\underline{\omega}_{\calE_1/C}^{\oplus 6}$ respectively. Write $\pi: S = \bbP_{\Fq}^1 \times F_d \to F_d$ be the projection and let $\sigma: \widetilde{S} \to S$ be the double cover as in \Cref{context:quadratic_twist_of_a_family_of_elliptic_curves_by_a_universal_family_of_polynomials}. The pullback of $\calE_1 \to C$ via $\sigma \circ p$ is the total space $\sigma^* \calE_{\text{univ}} \to \widetilde{S}$, which has equation
    $$y^2 = x^3 + (\sigma \circ \pi)^* A \cdot x + (\sigma \circ \pi)^* B.$$
    The involution $(Q,w,t,c) \mapsto (-Q,-w,t,c)$ on $\sigma^* \calE_{\text{univ}}$ fixes $x$ and negates $y$ and $w$. Setting
    $$X = x, \quad Y = wy,$$
    these are invariant under the action since $(-w)(-y) = wy$. These satisfy
    $$Y^2 = w^2 y^2 = P(t,c) \cdot \left( X^3 + (\sigma \circ \pi)^* A \cdot X + (\sigma \circ \pi)^* B \right),$$
    which is an equation that descends to $S$ because $P(t,c)$, $(\sigma \circ \pi)^* A$, and $(\sigma \circ \pi)^* B$ do not depend on $w$. In other words, we have a family of elliptic curves over $S$ with Weierstrass equation given by
    $$Y^2 = P(t,c) \cdot (X^3 + p^* A \cdot X + p^* B).$$
    The fibers of this relative generalized elliptic curve over $(t_0, C) \in S(\Fq)$ is thus indeed the quadratic twist of $\calE_1$ by $P(t_0,c_0)$.
\end{proof}

\begin{proposition}
    Assume \Cref{context:quadratic_twist_of_a_family_of_elliptic_curves_by_a_universal_family_of_polynomials}. Let $\calG = \sigma^* \calE_{\text{univ}} / (\bbZ/2\bbZ)$, which is a family of generalized elliptic curves over $S$. 
    \TODO{semistable}
    Let $j: S_{\text{sm}} \hookrightarrow S$ be the open immersion of the semistable locus, i.e. the locus where $\calG \to S$ has good or multiplicative reduction. The kernel $\calG[\ell]$ of multiplication by $\ell$ is an \'etale group scheme over $S$ that restricts to a finite group scheme over $S_{\text{sm}}$. We identify $\calG[\ell]$ with a an \'etale $\bbF_\ell$-sheaf $\calG_\ell$ on $S$ whose restriction to $S_{\text{sm}}$ is lisse. In particular, $j_* j^* \calG_\ell$ is a constructible $\bbF_\ell$-sheaf on $S$. Let $\pi: S = \bbP_{\Fq}^1 \times F_d \to F_d$ be the projection. The first higher direct image
    $$\hlin{\calT_{d,\ell} = R^1 \pi_* (j_* j^* \calG_\ell)}$$
    is a lisse $\bbF_\ell$-sheaf on $F_d$ and for every geometric point $\overline{f} \in F_d(\Fqbar)$, the geometric fiber is 
    $$\calT_{d,\ell}|_{\overline{f}} \cong H^1(\bbP_{\Fqbar}^1, \calE_{f,\ell})$$
    where \hl{$\calE_{f,\ell}$} is the N\'eron model of $\ell$-torsion subgroup $E_f[\ell]$ for the quadratic twist $E_f/\Fq(t)$.
\end{proposition}
\begin{proof}
    We have a Cartesian diagram:
    \begin{center}
    \begin{tikzcd}
    \mathbb{P}^1_{\overline{\mathbb{F}}_q} \arrow[r, "i_{\overline{f}}"] \arrow[d, "\pi_{\overline{f}}"] & S = \mathbb{P}^1_{\mathbb{F}_q} \times F_d \arrow[d, "\pi"] \\
    \text{Spec } \overline{\mathbb{F}}_q \arrow[r, "\overline{f}"] & F_d
    \end{tikzcd}
    \end{center}
    Since $\pi: S \to F_d$ is a proper morphism and $j_* j^* \calG_\ell$ is a constructible $\mathbb{F}_\ell$-sheaf on $S$, the Proper Base Change Theorem yields a canonical isomorphism:
    \TODO{proper base change theorem in this context}
    \[ \calT_{d,\ell}|_{\overline{f}} = (R^1 \pi_* (j_* j^* \calG_\ell))_{\overline{f}} \cong H^1(\mathbb{P}^1_{\overline{\mathbb{F}}_q}, i_{\overline{f}}^* j_* j^* \calG_\ell). \]

    Let $j_{\overline{f}}: S_{\text{sm}} \times_S \mathbb{P}^1_{\overline{\mathbb{F}}_q} \hookrightarrow \mathbb{P}^1_{\overline{\mathbb{F}}_q}$ denote the pullback of the open immersion $j$ via the inclusion $i_{\overline{f}}$:
    \begin{center}
    \begin{tikzcd}
        S_{\text{sm}} \times_S \mathbb{P}^1_{\overline{\mathbb{F}}_q} \arrow[r, "i_{\overline{f}, \text{sm}}"] \arrow[d, "j_{\overline{f}}"] & S_{\text{sm}} \arrow[d, "j"] \\
        \mathbb{P}^1_{\overline{\mathbb{F}}_q} \arrow[r, "i_{\overline{f}}"] & S.
    \end{tikzcd}
    \end{center}
    By the base change theorem for direct images of unramified torsion sheaves along open immersions, we have a canonical isomorphism:
    \[ i_{\overline{f}}^* j_* j^* \calG_\ell \cong (j_{\overline{f}})_* i_{\overline{f}, \text{sm}}^* j^* \calG_\ell. \]

    By \Cref{proposition:fiber_of_quadratic_twist_of_family_of_neron_model_of_elliptic_curves_over_F_q_t_over_polynomials_at_a_polynomial_is_quadaratic_twist_of_elliptic_curve_by_polynomial}, the restriction of $\calG$ to the fiber $\mathbb{P}^1_{\overline{\mathbb{F}}_q}$ is canonically isomorphic to the N\'eron model $\calE_f$. Since the formation of $\ell$-torsion kernels commutes with base change, we have the following sequence of canonical isomorphisms of group schemes over the smooth locus of the fiber:
    \[
    i_{\overline{f}, \text{sm}}^* \left( \calG|_{S_{\text{sm}}}[\ell] \right) \cong \left( i_{\overline{f}, \text{sm}}^* \calG|_{S_{\text{sm}}} \right)[\ell] \cong \left( j_{\overline{f}}^* \calE_f^{\text{sm}} \right)[\ell] \cong j_{\overline{f}}^* \left( \calE_f^{\text{sm}}[\ell] \right).
    \]
    Passing to the associated \'etale $\mathbb{F}_\ell$-sheaves yields the identification $i_{\overline{f}, \text{sm}}^* j^* \calG_\ell \cong j_{\overline{f}}^* \calE_{f,\ell}$. Substituting this into the above base change isomorphism gives:
    \[ i_{\overline{f}}^* j_* j^* \calG_\ell \cong (j_{\overline{f}})_* j_{\overline{f}}^* \calE_{f,\ell}. \]

    Since $\ell \nmid q$, the $\ell$-torsion subgroup scheme of a N\'eron model satisfies the extension property $\calE_{f,\ell} \cong (j_{\overline{f}})_* j_{\overline{f}}^* \calE_{f,\ell}$. We thereby obtain the desired identification of the geometric fiber:
    \[ \calT_{d,\ell}|_{\overline{f}} \cong H^1(\mathbb{P}^1_{\overline{\mathbb{F}}_q}, \calE_{f,\ell}). \]
    
    \TODO{verify the below}
    It remains to show that $\calT_{d,\ell}$ is a lisse $\mathbb{F}_\ell$-sheaf on $F_d$. Since $F_d$ is smooth over $\mathbb{F}_q$, it suffices to show that the dimension of the geometric fibers $\dim_{\mathbb{F}_\ell} H^1(\mathbb{P}^1_{\overline{\mathbb{F}}_q}, \calE_{f,\ell})$ is constant for all $\overline{f} \in F_d(\overline{\mathbb{F}}_q)$ and that the sheaf is unramified. 
    
    The dimension of $H^1(\mathbb{P}^1_{\overline{\mathbb{F}}_q}, \calE_{f,\ell})$ is given by the Grothendieck--Ogg--Shafarevich formula applied to the elliptic curve $E_f / \overline{\mathbb{F}}_q(t)$:
    \[ \dim_{\mathbb{F}_\ell} H^1(\mathbb{P}^1_{\overline{\mathbb{F}}_q}, \calE_{f,\ell}) = 2 \cdot \text{deg}(\text{Cond}(E_f)) - 4 + \sum_{v \in \mathbb{P}^1} \delta_v(E_f[\ell]). \]
    Because $\ell \neq p$, the wild ramification terms vanish ($\delta_v = 0$). Furthermore, since $f$ is a square-free polynomial of degree $d$ coprime to the conductor of $E_1$, the degree of the conductor divisor $\text{deg}(\text{Cond}(E_f))$ is constant across all $\overline{f} \in F_d(\overline{\mathbb{F}}_q)$. Thus, the higher direct image $\calT_{d,\ell}$ has constant fiber dimension and is locally constant, hence it is a lisse $\mathbb{F}_\ell$-sheaf on $F_d$.
    
    % By \Cref{prop:universal_quadratic_twist}, the restriction of $\calG$ to the fiber over $\overline{f}$ is canonically isomorphic to the smooth locus of the N\'eron model $\calE_f \to \mathbb{P}^1_{\overline{\mathbb{F}}_q}$. Taking $\ell$-torsion preserves this isomorphism on the smooth loci, so $i_{\overline{f}, \text{sm}}^* j^* \calG_\ell \cong j_{\overline{f}}^* \calE_{f,\ell}$. Substituting this back yields:
    % \[ i_{\overline{f}}^* j_* j^* \calG_\ell \cong (j_{\overline{f}})_* j_{\overline{f}}^* \calE_{f,\ell}. \]
    
    % By the extension property of the $\ell$-torsion subgroup scheme of a N\'eron model when $\ell \nmid q$, the natural map $\calE_{f,\ell} \to (j_{\overline{f}})_* j_{\overline{f}}^* \calE_{f,\ell}$ is an isomorphism. Consequently, we obtain the desired identification of the geometric fiber:
    % \[ \calT_{d,\ell}|_{\overline{f}} \cong H^1(\mathbb{P}^1_{\overline{\mathbb{F}}_q}, \calE_{f,\ell}). \]
    
    % It remains to show that $\calT_{d,\ell}$ is a lisse $\mathbb{F}_\ell$-sheaf on $F_d$. Since $F_d$ is smooth over $\mathbb{F}_q$, it suffices to show that the dimension of the geometric fibers $\dim_{\mathbb{F}_\ell} H^1(\mathbb{P}^1_{\overline{\mathbb{F}}_q}, \calE_{f,\ell})$ is constant for all $\overline{f} \in F_d(\overline{\mathbb{F}}_q)$ and that the sheaf is unramified. 
    
    % The dimension of $H^1(\mathbb{P}^1_{\overline{\mathbb{F}}_q}, \calE_{f,\ell})$ can be computed via the Grothendieck--Ogg--Shafarevich formula applied to the elliptic curve $E_f$ over $\overline{\mathbb{F}}_q(t)$:
    % \[ \dim_{\mathbb{F}_\ell} H^1(\mathbb{P}^1_{\overline{\mathbb{F}}_q}, \calE_{f,\ell}) = 2 \cdot \text{deg}(\text{Cond}(E_f)) - 4 + \sum_{v \in \mathbb{P}^1} \delta_v(E_f[\ell]) \]
    % where $\text{Cond}(E_f)$ is the conductor divisor and $\delta_v$ denotes the local wild ramification terms. Since $\ell \neq p$, the wild ramification terms vanish ($\delta_v = 0$). 
    
    % Furthermore, because $f$ is a square-free polynomial of degree $d$ coprime to the conductor of $E_1$, the conductor of the quadratic twist $E_f$ changes in a totally predictable, uniform manner based purely on the degree $d$ and the bad reduction locus of $E_1$. Thus, the degree of the conductor $\text{deg}(\text{Cond}(E_f))$ is completely invariant across all $\overline{f} \in F_d(\overline{\mathbb{F}}_q)$. 
    
    % Because the dimensions of the fibers are constant and the underlying sheaf sequence transitions smoothly in a family over the smooth base parameter space $F_d$, the higher direct image $\calT_{d,\ell}$ is locally constant, and hence a lisse $\mathbb{F}_\ell$-sheaf on $F_d$.
\end{proof}

\begin{proposition} \label{proposition:geometric_stalk_of_pushforward_of_lisse_sheaf_on_open_subscheme_of_integral_normal_scheme_at_codimension_1_point_is_invariant_subgroup_of_representation_by_inertia_group}
    Let $X$ be an \CrefAndHyperrefIfExist{definition:integral_scheme}{integral} \CrefAndHyperrefIfExist{definition:normal_and_factorial_scheme}{normal scheme} with \CrefAndHyperrefIfExist{definition:generic_point_of_a_topological_space}{generic point} $\eta$. Let $j: U \hookrightarrow X$ be the inclusion of a non-empty \CrefAndHyperrefIfExist{definition:open_and_closed_subscheme_of_a_scheme}{open subscheme}. Let $x \in X \setminus U$ be a point of \CrefAndHyperrefIfExist{definition:codimension_of_an_irreducible_closed_subset_in_a_topological_space}{codimension} $1$ in $X$. Write $I_x \subset \pioneet(U,\overline{\eta})$ for the \CrefAndHyperrefIfExist{definition:decomposition_and_inertia_subgroups_of_etale_fundamental_group_of_dense_open_of_integral_normal_scheme_at_point_in_complement}{inertia group}. 

    % Let $C$ be a \CrefAndHyperrefIfExist{definition:smooth_morphism_of_schemes}{smooth}, \CrefAndHyperrefIfExist{definition:connected_object_of_a_galois_category}{connected}, \CrefAndHyperrefIfExist{definition:algebraic_curve_over_a_field}{algebraic curve} over an \CrefAndHyperrefIfExist{definition:algebraically_closed_field}{algebraically closed} \CrefAndHyperrefIfExist{definition:field}{field} $k$. Let $S \subset X$ be a \CrefAndHyperrefIfExist{definition:countable_finite_uncountable_sets}{finite set} of closed points, and let $U = X \setminus S$ be the open complement. Let $j: U \hookrightarrow X$ be the \CrefAndHyperrefIfExist{definition:open_and_closed_immersions_of_schemes}{open immersion}. For any \CrefAndHyperrefIfExist{definition:closed_point_of_a_topological_space}{closed point} $x \in S$, let $\overline{x} \to X$ be a chosen \CrefAndHyperrefIfExist{definition:geometric_point_of_a_scheme}{geometric point} over $x$. This choice determines a local \CrefAndHyperrefIfExist{definition:decomposition_and_inertia_subgroups_of_etale_fundamental_group_of_dense_open_of_integral_normal_scheme_at_point_in_complement}{decomposition group} at $x$, which canonically identifies with the local \CrefAndHyperrefIfExist{definition:decomposition_and_inertia_subgroups_of_etale_fundamental_group_of_dense_open_of_integral_normal_scheme_at_point_in_complement}{inertia group} $I_x \subset \pioneet(U,\overline{\eta})$ where $\eta$ is the generic point of $U$. 

    Let $\calG$ be a \CrefAndHyperrefIfExist{definition:lisse_sheaf_on_a_scheme}{lisse sheaf} on $U$, in any of the following sense:
    \begin{enumerate}
        \item a $\Lambda$-lisse sheaf where $\Lambda$ is a finite ring,
        \item An adic lisse sheaf with coefficients in $R$, $E$, or $K$, where $(R,\lambda)$ is a complete DVR with $R/\lambda^n$ finite for all $n \geq 1$, $E$ is its fraction field, and $K$ is an algebraic extension of $E$, 
    \end{enumerate}
    Write $\rho: \pioneet(U,\overline{\eta}) \to \GL(V)$ for the corresponding\CrefIfExists{theorem:locally_constant_sheaves_on_X_etale_are_equivalent_to_representations_of_the_etale_fundamental_group}\CrefIfExists{theorem:lisse_ell_adic_sheaves_on_connected_noetherian_schemes_are_equivalent_to_representations_of_etale_fundamental_group_over_adic_ring} continuous representation where $V$ is a module over an appropriate ring.

    The geometric stalk $(j_* \calG)_{\overline{x}}$
    \CrefIfExists{definition:stalk_of_a_presheaf_on_the_small_etale_site_of_a_schem_at_a_geometric_point}\CrefIfExists{definition:geometric_stalk_of_an_etale_lambda_adic_sheaf_on_a_scheme}\CrefIfExists{definition:geometric_stalk_of_an_etale_lambda_adic_sheaf_on_a_scheme_with_rational_coefficients} is isomorphic to $V^{I_{\overline{x}}}$ \CrefIfExists{definition:invariant_subgroup_of_a_group_acting_on_an_abelian_group}.
\end{proposition}
\begin{proof}
    It suffices to prove this in the case that $\calG$ is a lisse sheaf with coefficients in $\Lambda$ for a finite ring $\Lambda$. By \Cref{theorem:stalks_of_etale_sheaves_at_geometric_points_are_the_values_of_the_sheaves_at_the_strict_henselization_of_the_local_rings_at_the_image_of_the_point}, we have 
    $$(j_* \calG)_{\overline{x}} \cong \Gamma(\Spec(\calO_{X,x}^{\text{sh}}), g^* j_* \calG)$$
    where $g: \Spec(\calO_{X,x}^{\text{sh}}) \to X$ is the canonical morphism. Let $K_x^{\text{sh}}$ be the fraction field of $\calO_{X,x}^{\text{sh}}$. We have the Cartesian square
    \begin{center}
    \begin{tikzcd}
        K_x^{\text{sh}} \ar[r, "\widetilde{g}"]   \ar[d, "\widetilde{j}"] & U \ar[d, "j"] \\ 
        \Spec(\calO_{X,x}^{\text{sh}}) \ar[r, "g"] & X.  
    \end{tikzcd}
    \end{center}
    By the base change theorem, \TODO{which one}
    $$g^* j_* \calG \cong \widetilde{j}_*\widetilde{g}^* \calG$$
    and hence
    $$(j_* \calG)_{\overline{x}} \cong \Gamma(\Spec(\calO_{X,x}^{\text{sh}}), \widetilde{j}_*\widetilde{g}^* \calG) \cong \Gamma(\Spec(K_x^{\text{sh}}) , \widetilde{g}^* \calG)$$
    \TODO{write as a general fact that inertia group inclusion corresponds to scheme map $\Spec(K_x^{\text{sh}}) \to U$}
    The inclusion $\widetilde{g}: \Spec(K_x^{\text{sh}}) \to U$ induces the inclusion $I_{x} \hookrightarrow \pioneet(U, \overline{\eta})$ of the inertia group at $x$. In particular, the pullback $\widetilde{g}^* \calG$ corresponds to the representation $\rho$ restricted to $I_{x}$ and its global sections is the $I_x$-invariants of $V$. In other words, $(j_* \calG)_{\overline{x}} \cong V^{I_x}$ as desired.
\end{proof}

% \begin{proof}
%     We have a Cartesian diagram
%     \begin{center}
%     \begin{tikzcd}
%     \mathbb{P}^1_{\overline{\mathbb{F}}q} \arrow[r, "i_{\overline{f}}"] \arrow[d, "\pi_{\overline{f}}"] & S = \mathbb{P}^1_{\mathbb{F}_q} \times F_d \arrow[d, "\pi"] \\
%     \text{Spec } \overline{\mathbb{F}}_q \arrow[r, "\overline{f}"] & F_d
%     \end{tikzcd}
%     \end{center}
%     By the proper base change theorem,
%     $$(R^1 \pi_* (j_* j^* \calG_\ell))_{\overline{f}} \cong H^1(\bbP_{\Fqbar}^1, i_{\overline{f}}^* j_* j^* \calG_\ell )$$

%     \begin{center}
%     \begin{tikzcd}
%         S_{\text{sm}} \cap (\bbP_{\Fqbar}^1)  \ar[r] \ar[d]  & \Spec \Fqbar \ar[d, "\overline{f}"]  \\
%     S_{\text{sm}} \ar[r, "j \circ \pi"]  & F_d 
%     \end{tikzcd}
%     \end{center}
%     Write $\Spec \Fqbar \to $
% \end{proof}

\input{../_assembly/source_bridge/theorem_explicit_construction_of_T_d_ell_as_derived_pushforward_of_Neron_l_torsion_sheaf_on_universal_quadratic_cover.tex}


\section{Ellenberg and Landesman's Needs}

\begin{definition}[see {\cite[Definition 2.4.2.]{ellenberg_landesman_hsghssgqtfff}} for a discussion] \label{definition:hurwitz_stack_of_G_covers_over_a_curve_over_a_base_scheme}
    \TODO{divisor over a scheme}
    Let $B$ be a \CrefAndHyperrefIfExist{definition:scheme}{base scheme}. Let $C \to B$ be a \CrefAndHyperrefIfExist{definition:smooth_morphism_of_schemes}{smooth}, \CrefAndHyperrefIfExist{definition:proper_morphism_of_schemes}{proper} \CrefAndHyperrefIfExist{definition:curve_over_a_scheme}{relative curve} with \CrefAndHyperrefIfExist{definition:has_geometrically_connected_fibers_for_a_morphism_of_schemes}{geometrically connected fibers} of \CrefAndHyperrefIfExist{definition:geometric_genus_of_a_smooth_projective_variety_over_a_field}{genus} $g$. Let $Z \subset C$ be a \CrefAndHyperrefIfExist{definition:finite_etale_cover_of_a_scheme}{finite \'etale} \CrefAndHyperrefIfExist{definition:relative_effective_cartier_divisor_on_a_scheme_over_a_scheme}{divisor} over $B$. Let $U = C \setminus Z$. 

    \TODO{boundary compnent, puncture}
    Let $\Sigma_{g,n,f+1}$ a \CrefAndHyperrefIfExist{definition:genus_of_a_topological_space}{genus} $g$ \CrefAndHyperrefIfExist{definition:surface_as_a_manifold}{topological surface} with $b$ boundary components and $f$ punctures.
    \TODO{action}
    Suppose that $\mathcal{S} \subset \Hom(\pioneet(\Sigma_{g,n,f+1}),G)$ is a $G$ conjugation invariant subset preserved by the action of $\pioneet(\Conf_{U_{\overline{b}} / \overline{b}}^n)$ acting on the first $n$ points. Define the \hldef{Hurwitz stack} $\Hur_{C/B}^{G,n,Z,S}$ to be the \CrefAndHyperrefIfExist{definition:artin_stack_over_a_scheme_for_a_grothendieck_topology_on_the_category_of_S_schemes}{algebraic stack} over $B$ whose functor of points is defined as follows: for $T$ a $B$-scheme, $\Hur_{C/B}^{G,n,Z,S}(T)$ is the \CrefAndHyperrefIfExist{definition:groupoid}{groupoid} whose objects are the tuples
    $$(D,i: D \to C_T, X, h: X \to C_T)$$
    such that 
    \begin{itemize}
        \item $D$ is a \CrefAndHyperrefIfExist{definition:finite_etale_cover_of_a_scheme}{finite \'etale cover} of $T$ of degree $n$.
        \item $i$ is a \CrefAndHyperrefIfExist{definition:open_and_closed_immersions_of_schemes}{closed immersion} $i: D \subset C_T$ that is disjoint from $Z_T \subset C_T$.
        \item $X$ is a smooth proper relative curve over $T$, not necessarily having geometrically connected fibers.
        \item $h: X \to C_T$ is a \CrefAndHyperrefIfExist{definition:finite_locally_free_morphism_of_schemes}{finite locally free} \CrefAndHyperrefIfExist{definition:galois_morphism_of_schemes}{Galois $G$-cover} that is \'etale away from $Z_T \cup i(D) \subset C_T$. 
        \item Let ${\overline{t}} \to T$ be a fixed geometric point. Let $\overline{\eta}$ denote the geometric generic point of $(C_T)_{\overline{t}}$. Then the representation $\rho: \pi_1((U_T)_{\overline{t}}-i(D_{\overline{t}}), \overline{\eta}) \to G$ \CrefAndHyperrefIfExist{theorem:finite_etale_covers_of_a_connected_scheme_correspond_to_finite_sets_with_an_action_of_the_etale_fundamental_group}{corresponding to} the $G$-cover $h$, under the identification of $\mathrm{Hom}(\pi_1((U_T)_{\overline{t}}-i(D_{\overline{t}}), \overline{\eta}),G)$ and $\mathrm{Hom}(\pi_1(\Sigma_{g,n+f+1}), G)$ corresponds to an element of $\mathcal S$.
    \end{itemize}
    Two covers $X \to C_T$ are considered equivalent if they are related by the $G$-conjugation action \TODO{what does this mean}. Morphisms between two points $(D_j, i_j, X_j, h_j)$ for $j \in \{1,2\}$ are given by $(\phi_D, \psi_X)$ where $\phi_D: D_1 \simeq D_2$ is an isomorphism so that $i_2 \circ \phi_D = i_2$ and $\psi_X: X_1 \simeq X_2$ is an isomorphism such that $h_2 \circ \psi_X = h_1$ and $\psi_X = g^{-1} \psi_X g$ for every $g \in G$.
\end{definition}

\begin{definition}
    Let $B$ be a \CrefAndHyperrefIfExist{definition:scheme}{base scheme}. Let $C \to B$ be a \CrefAndHyperrefIfExist{definition:smooth_morphism_of_schemes}{smooth}, \CrefAndHyperrefIfExist{definition:proper_morphism_of_schemes}{proper} \CrefAndHyperrefIfExist{definition:curve_over_a_scheme}{relative curve} with \CrefAndHyperrefIfExist{definition:has_geometrically_connected_fibers_for_a_morphism_of_schemes}{geometrically connected fibers} of \CrefAndHyperrefIfExist{definition:geometric_genus_of_a_smooth_projective_variety_over_a_field}{genus} $g$. Let $Z \subset C$ be a \CrefAndHyperrefIfExist{definition:finite_etale_cover_of_a_scheme}{finite \'etale} \CrefAndHyperrefIfExist{definition:relative_effective_cartier_divisor_on_a_scheme_over_a_scheme}{divisor} over $B$. Let $U = C \setminus Z$. 

    
\end{definition}

\begin{proposition}
    Let $B = \overline{b} = \Spec k$ be a geometric point where $k$ is an algebraically closed field.
    
    \TODO{say what g,r,f, Dx, F} are
    Let $n$ be an even integer such that 
    \begin{align*}
        n > \max\left(2g,\frac{2(2r+1)(f+1) - \sum_{y \in D_x({\overline b})} \mathrm{Drop}_y(\mathscr F)}{2r} - (2g-2) \right).
    \end{align*}
    \CrefIfExists{definition:drop_of_a_lisse_sheaf_on_dense_open_of_integral_normal_scheme_at_point_in_complement}
    Let $C \to B$ be a smooth, proper connected curve of genus $g$. Let $Z \subset C$ be a finite \'etale divisor over $B$. Let $U = C \setminus Z$. 
    Let $\operatorname{QTwist}_{U/B}^n \coloneq \operatorname{Hur}_{C/B}^{\bbZ/2\bbZ, n, Z, \mathcal{S}}$(\Cref{definition:hurwitz_stack_of_G_covers_over_a_curve_over_a_base_scheme}); note that there is a natural projection map 
    $$\operatorname{QTwist}_{U/B}^n \to \Conf_{U/B}^n.$$
    \TODO{this variant of the configuration space}
    Let $x = [(D,\phi)] \in \operatorname{QTwist}_{U/B}^n$ be a $k$-point, let $y \in \Conf_{U/B}^n$ be its image under the above projection map. Let 
    $$\mathscr{C}_B^n \coloneq \Conf_{U/B}^n \times_B C$$
    be the ``universal curve'' over the configuration space $\Conf_{U/B}^n$ associated to the curve $C \to B$, and let 
    $$\xi: \mathscr{C}_B^n \to \Conf_{U/B}^n$$
    be the projection map. Let $C_x \cong y \times_B C \cong C$ be the fiber of $\xi$ over $y$. 

    By \TODO{cite katz}, there is a map $h: C_x \to \mathbb{P}^1$ of degree $n$ such that the branch locus $\operatorname{br}(h)$ of $h$ is disjoint from $h((Z \cup D)_x)$, $h$ separates points of $(Z \cup D)_x$, and precisely one point $\delta \in D_x$ maps to $\infty \in \bbP^1$. Let $W \subset \mathbb{P}^1$ to be the complement of $\operatorname{br}(h) \cup h(Z \cup D)$. 

    \TODO{wait, how do we define kummer sheaves over more general fields?}
    We construct a map $\phi: W \to \operatorname{QTwist}_{U/\overline{b}}^n$. Such a map corresponds to a finite double cover $V$ of $C \times_{\overline{b}} W$, which in turn is equivalent to specifying a rank $1$ locally constant constructible $\bbZ/\ell\bbZ$-sheaf on an open subset of $C \times_{\overline{b}} W$ whose monodromy is trivialized by. Let $t: X \to U$ be the finite double cover corresponding to $x$; in particular, $t$ is \'etale away from $D$. Let $\bbV \coloneq t_* (\bbZ/\ell \bbZ)/(\bbZ/\ell\bbZ)$, which is a rank $1$ constructible sheaf $\bbZ/\ell \bbZ$-sheaf on $U$ that becomes locally constant on $U-D$. Let $\mathscr{F}' = \mathscr{F} \otimes \mathbb{V}$, i.e. $\mathscr{F}'$ is the quadratic twist of $\mathscr{F}$ corresponding to $x$.  \TODO{think about how $\mathscr{F}'$ is the quadratic twist corresponding to $x$}
    
    Let $\calL_\chi$ be the \CrefAndHyperrefIfExist{definition:kummer_sheaf_associated_to_a_character_on_points_of_a_connected_commutative_algebraic_group_over_a_finite_field}{Kummer sheaf} on $\bbG_m$ corresponding to the double cover $\bbG_m \xrightarrow{z \mapsto z^2} \bbG_m$. Construct the following maps:
    \begin{align*}
    \alpha'&: \bbA^1 \times \bbA^1 -\Delta \to \bbG_m, \quad (x,y) \mapsto x-y \\
    \alpha&: Y \coloneq (h,\id)^{-1}(W \times W - \Delta) \xrightarrow{(h,\id)} \bbA^1 \times \bbA^1 - \Delta \xrightarrow{\alpha'} \bbG_m.
    \end{align*}
    let $\mathbb{W} \coloneq \alpha^* \calL_\chi$. 
    Let $p_C: C \times W \to C$ be the projection.
    % Let $\pi_2: Y \to \bbA^1$ be the second projection. 
    \TODO{continue} 
    Noting that $h^{-1}(W-0) \subset U-D$, let $\mathbb{V}'$ be the sheaf on $h^{-1}(W-0)$ given by 
    $$\mathbb{V}' \coloneq \mathbb{W}|_{h^{-1}(W-0) \times 0} \otimes \mathbb{V}|_{h^{-1}(W-0)}.$$
    

\end{proposition}

\section{Twists}

\begin{definition} \label{definition:twist_of_an_etale_sheaf_on_a_scheme_with_various_coefficients_by_a_character_on_the_etale_fundamental_group}
    \TODO{AI generated, verify}
    Let $U$ be a \CrefAndHyperrefIfExist{definition:disconnected_and_connected_topological_spaces}{connected} \CrefAndHyperrefIfExist{definition:scheme}{scheme}, and let $\bar{u}: \Spec(\Omega) \to U$ be a \CrefAndHyperrefIfExist{definition:geometric_point_of_a_scheme}{geometric point}. Let $\pioneet(U, \bar{u})$ denote the \CrefAndHyperrefIfExist{definition:etale_fundamental_group_of_a_connected_scheme}{\'etale fundamental group} of $U$ at $\bar{u}$.

    Let $\Lambda$ be one of the following coefficient rings:
    \begin{enumerate}
        \item A finite ring,
        \item A complete DVR $R$,
        \item The \CrefAndHyperrefIfExist{definition:field_of_fractions_of_an_integral_domain}{field of fractions} $E$ of a \CrefAndHyperrefIfExist{definition:complete_ring_with_respect_to_an_ideal}{complete} DVR $R$, or
        \item An \CrefAndHyperrefIfExist{definition:algebraic_element_minimal_polynomial_of_an_algebraic_element_algebraic_field_extension_transcendental_element_field_extension}{algebraic extension} of the \CrefAndHyperrefIfExist{definition:field_of_fractions_of_an_integral_domain}{field of fractions} $E$ of a complete DVR $R$.
    \end{enumerate}
    If $\Lambda$ is a complete DVR, a fraction field thereof, or an algebraic extension of such a fraction field, then further assume that $U$ is a \CrefAndHyperrefIfExist{definition:disconnected_and_connected_topological_spaces}{connected} \CrefAndHyperrefIfExist{definition:locally_noetherian_and_noetherian_scheme}{noetherian scheme}.

    \begin{enumerate}
        \item Let $A$ be one of the coefficient rings listed above. A \hldef{character on $\pioneet(U, \bar{u})$ with values in $A$} is a continuous group homomorphism
        $$\hlin{\chi: \pioneet(U, \bar{u}) \to A^\times}$$
        where $A^\times$ denotes the group of units of $A$, equipped with the discrete topology.

        \item Let $\chi: \pioneet(U, \bar{u}) \to A^\times$ be a character as in (1). By the equivalence of categories between lisse sheaves on $U_{\text{'\'et}}$ and continuous representations of $\pioneet(U, \bar{u})$
        \CrefIfExists{theorem:locally_constant_sheaves_on_X_etale_are_equivalent_to_representations_of_the_etale_fundamental_group}\CrefIfExists{theorem:lisse_ell_adic_sheaves_on_connected_noetherian_schemes_are_equivalent_to_representations_of_etale_fundamental_group_over_adic_ring}, the character $\chi$ corresponds uniquely to a rank-1 \CrefAndHyperrefIfExist{definition:lisse_sheaf_on_a_scheme}{lisse sheaf}, denoted $\mathcal{L}_\chi$, often called the \hldef{rank-1 lisse sheaf associated to $\chi$}. The stalks of $\mathcal{L}_\chi$ are isomorphic to $A$ as $A$-modules, and the action of $\gamma \in \pioneet(U, \bar{u})$ on a stalk is multiplication by $\chi(\gamma)$.

        \item Let $\mathscr{F}$ be a \CrefAndHyperrefIfExist{definition:lisse_sheaf_on_a_scheme}{lisse sheaf} with coefficients in $A$, corresponding to a representation $\rho: \pioneet(U, \bar{u}) \to \text{Aut}_A(V)$, and let $\chi$ be a character as above. The \hldef{twist of $\mathscr{F}$ by $\chi$}, denoted \hl{$\mathscr{F}_\chi$}, is the lisse sheaf $\mathscr{F} \otimes \mathcal{L}_\chi$\CrefIfExists{definition:sheafified_tensor_product_of_sheaves_of_modules_of_sheaves_of_rings_on_a_site_with_a_small_topologically_generating_family} corresponding to the tensor product representation $$(\rho \otimes \chi)(\gamma) = \chi(\gamma) \cdot \rho(\gamma)$$
        for all $\gamma \in \pioneet(U, \bar{u})$. Since $\mathcal{L}_\chi$ has rank 1, the stalks of $\mathscr{F} \otimes \mathcal{L}_\chi$ are isomorphic to those of $\mathscr{F}$ as $A$-modules, but the monodromy action is scaled by the character $\chi$.

        \item Two lisse sheaves $\mathscr{F}$ and $\mathscr{G}$ with coefficients in $A$ are called \hldef{twists of each other} if there exists a character $\chi: \pioneet(U, \bar{u}) \to A^\times$ such that
        $$\mathscr{G} \cong \mathscr{F} \otimes \mathcal{L}_\chi.$$
    \end{enumerate}
\end{definition}



\appendix


\section{Abelian varieties}
\input{../_definitions/definition_algebraic_variety_over_a_field.tex}
\input{../_definitions/definition_algebraic_group_scheme_over_a_scheme.tex}
\input{../_definitions/definition_homomorphism_of_algebraic_groups_over_a_scheme.tex}
\begin{definition}[Kernel of a homomorphism of group schemes over a scheme] \label{definition:kernel_of_a_homomorphism_of_group_schemes_over_a_scheme}
    Let $S$ be a \CrefAndHyperrefIfExist{definition:scheme}{scheme} and let $\phi\colon G_1 \to G_2$ be a \CrefAndHyperrefIfExist{definition:homomorphism_of_algebraic_groups_over_a_scheme}{homomorphism of algebraic group schemes over $S$}.

    The \hldef{kernel of $\phi$}, denoted by \hl{$\ker \phi$} or \hl{$A[\phi]$}, is the \CrefAndHyperrefIfExist{definition:fiber_of_a_morphism_of_schemes_at_a_point}{fiber} of $\phi$ at the \CrefAndHyperrefIfExist{definition:algebraic_group_scheme_over_a_scheme}{identity section} $e_2\colon S \to G_2$ of $G_2$. Equivalently, $\ker \phi$ is the subgroup scheme of $G_1$ defined as the \CrefAndHyperrefIfExist{definition:cartesian_product_of_two_objects_in_a_category_over_an_object}{pullback}:
    $$\hlin{\ker \phi = G_1 \times_{G_2} S,}$$
    \TODO{finite group scheme}
    where the map $G_1 \to G_2$ is $\phi$ and the map $S \to G_2$ is the identity section. If $\phi$ is a \CrefAndHyperrefIfExist{definition:finite_morphism_of_schemes}{finite morphism}, then $\ker \phi$ is a finite group scheme over $S$.
\end{definition}

\TODO{group variety}
\input{../_definitions/definition_abelian_variety_over_a_field.tex}
\begin{definition}[Isogeny of abelian varieties over a field] \label{definition:isogeny_of_abelian_varieties_over_a_field}
    \TODO{finite group scheme}
    Let $k$ be a \CrefAndHyperrefIfExist{definition:field}{field} and let $A$ and $B$ be \CrefAndHyperrefIfExist{definition:abelian_variety_over_a_field}{abelian varieties} over $k$. An \hldef{isogeny of abelian varieties over $k$} is a \CrefAndHyperrefIfExist{definition:homomorphism_of_algebraic_groups_over_a_scheme}{homomorphism} $\varphi\colon A \to B$ of group schemes over $k$ that is \CrefAndHyperrefIfExist{definition:injective_surjective_bijective_map_of_sets}{surjective} and has finite \CrefAndHyperrefIfExist{definition:kernel_of_a_homomorphism_of_group_schemes_over_a_scheme}{kernel}. The kernel $\ker \varphi$ is a finite \CrefAndHyperrefIfExist{definition:kernel_of_a_homomorphism_of_group_schemes_over_a_scheme}{group scheme over $k$}. Equivalently, $\varphi$ is a \CrefAndHyperrefIfExist{definition:finite_morphism_of_schemes}{finite}, \CrefAndHyperrefIfExist{definition:faithfully_flat_morphism_of_schemes}{faithfully flat} homomorphism of abelian varieties over $k$.
\end{definition}


\subsection{Picard groups and Jacobians}

\begin{definition} \label{definition:picard_group_of_a_scheme}
    \TODO{picard group of an algebraic space}
Let $X$ be a \CrefAndHyperrefIfExist{definition:scheme}{scheme}.
The \hldef{Picard group of $X$}, denoted by \hl{$\operatorname{Pic}(X)$}, is the \CrefAndHyperrefIfExist{definition:group}{group} of isomorphism classes of \CrefAndHyperrefIfExist{definition:free_and_locally_free_modules_on_a_ringed_site}{invertible sheaves} (line bundles) on $X$ (with respect to the \CrefAndHyperrefIfExist{definition:small_zariski_site_of_a_schem}{Zariski topology}), with the group operation induced by the \CrefAndHyperrefIfExist{definition:tensor_product_of_sheaves_of_bimodules_over_sheaves_of_rings_on_a_site}{tensor product of sheaves}.

More precisely:
\begin{enumerate}
    \item Two invertible \CrefAndHyperrefIfExist{definition:sheaf_on_a_site}{sheaves} $\mathcal{L}$ and $\mathcal{M}$ are \CrefAndHyperrefIfExist{definition:isomorphism_in_a_category}{isomorphic} if there exists an \CrefAndHyperrefIfExist{definition:module_over_a_sheaf_of_rings_on_a_ringed_space}{$\mathcal{O}_X$-module} isomorphism $\phi: \mathcal{L} \xrightarrow{\sim} \mathcal{M}$.
    \item The \CrefAndHyperrefIfExist{definition:tensor_product_of_sheaves_of_bimodules_over_sheaves_of_rings_on_a_site}{tensor product} of invertible sheaves induces a well-defined group operation on the set of isomorphism classes.
    \TODO{sheafy dual}
    \item The inverse of an \CrefAndHyperrefIfExist{definition:algebraic_vector_bundle_over_a_scheme}{invertible sheaf} $\mathcal{L}$ is given by its \CrefAndHyperrefIfExist{definition:dual_of_a_left_right_two_sided_module}{dual} $\mathcal{L}^\vee = \mathscr{H}om_{\mathcal{O}_X}(\mathcal{L}, \mathcal{O}_X)$\CrefIfExists{definition:sheaf_hom_for_sheaves_of_modules_over_sheaves_of_rings_on_a_site_with_a_small_topologically_generating_family}.
    \item The \CrefAndHyperrefIfExist{definition:left_right_identity_element_of_a_binary_operation_on_a_set}{identity element} is the class of the \CrefAndHyperrefIfExist{definition:ringed_space}{structure sheaf} \CrefAndHyperrefIfExist{definition:algebra_over_a_sheaf_of_rings_on_a_site}{$\mathcal{O}_X$}.
\end{enumerate}

\TextIfExists{definition:cartier_divisor_on_a_scheme}{
    There is a canonical isomorphism between $\operatorname{Pic}(X)$ and the \CrefAndHyperrefIfExist{definition:cartier_divisor_on_a_scheme}{Cartier divisor class group} $\operatorname{CaCl}(X)$ (\CrefAndHyperrefIfExist{definition:cartier_divisor_on_a_scheme}{Cartier divisor on a scheme}), given by associating to a Cartier \CrefAndHyperrefIfExist{definition:prime_and_weil_divisor_on_a_scheme}{divisor} $D$ the invertible sheaf $\mathcal{O}_X(D)$.
}

\TODO{class group of ring}
When $X = \operatorname{Spec}(A)$ is an \CrefAndHyperrefIfExist{definition:affine_scheme}{affine scheme}, there is a canonical isomorphism
$$\operatorname{Pic}(\operatorname{Spec}(A)) \cong \operatorname{Cl}(A),$$
where $\operatorname{Cl}(A)$ is the class group of the \CrefAndHyperrefIfExist{definition:ring}{ring} $A$ (the group of isomorphism classes of rank-1 \CrefAndHyperrefIfExist{definition:injective_and_projective_objects_in_a_category}{projective} modules over $A$, also known as the \CrefAndHyperrefIfExist{definition:picard_group}{Picard group} of the ring).

\TODO{extract this out}
If $X$ is a \CrefAndHyperrefIfExist{definition:noetherian_scheme}{Noetherian}, \CrefAndHyperrefIfExist{definition:integral_scheme}{integral}, and \CrefAndHyperrefIfExist{definition:locally_factorial_scheme}{locally factorial} \CrefAndHyperrefIfExist{definition:scheme}{scheme}, then $\operatorname{Pic}(X)$ is isomorphic to the \CrefAndHyperrefIfExist{definition:weil_divisor_class_group}{Weil divisor class group} $\operatorname{Cl}(X)$.
\end{definition}

\begin{definition} \label{definition:picard_variety_of_a_smooth_projective_geometrically_integral_variety_over_a_field}
    Let $X/k$ be a \CrefAndHyperrefIfExist{definition:smooth_at_a_point_for_a_finite_type_scheme_over_a_field}{smooth}, \CrefAndHyperrefIfExist{definition:geometric_property_of_a_scheme_over_a_field}{geometrically  integral} \CrefAndHyperrefIfExist{definition:projective_variety}{projective variety} over a \CrefAndHyperrefIfExist{definition:field}{field} $k$. 

    The \CrefAndHyperrefIfExist{definition:relative_picard_functor_of_a_scheme_over_a_scheme}{relative Picard functor} $\Pic_{X/k}$ is \CrefAndHyperrefIfExist{definition:representable_functor_on_a_locally_small_category}{representable by} a \CrefAndHyperrefIfExist{definition:smooth_morphism_of_schemes}{smooth scheme} over $k$. 
    \TODO{show that the identity component is an abelian variety}
    The \CrefAndHyperrefIfExist{definition:identity_component_of_a_group_scheme_over_a_scheme}{identity component} $\Pic_{X/k}^0$ is in fact an abelian variety and is called the \hldef{Picard variety of $X/k$}.
\end{definition}
\begin{definition}[Jacobian variety of a curve] \label{definition:jacobian_of_a_curve_over_a_field}

    \TODO{separate out jacobian of just a proper curve and the jacobian of a nice curve}
    Let $X$ be a \CrefAndHyperrefIfExist{definition:proper_morphism_of_schemes}{proper} \CrefAndHyperrefIfExist{definition:algebraic_curve_over_a_field}{curve} over a \CrefAndHyperrefIfExist{definition:field}{field} $k$.

    \begin{enumerate}
        \item The \CrefAndHyperrefIfExist{definition:relative_picard_functor_of_a_scheme_over_a_scheme}{relative Picard functor} $\Pic_{X/k}$ is \CrefAndHyperrefIfExist{definition:representable_functor_on_a_locally_small_category}{representable by} a \CrefAndHyperrefIfExist{definition:smooth_morphism_of_schemes}{smooth} \CrefAndHyperrefIfExist{definition:algebraic_group_scheme_over_a_scheme}{group scheme} over $k$. The \hldef{(generalized) Jacobian of $X$} refers to the \CrefAndHyperrefIfExist{definition:identity_component_of_a_group_scheme_over_a_scheme}{identity component} $\Pic_{X/k}^0$ of this group scheme.

        \item Now assume that $X$ is \CrefAndHyperrefIfExist{definition:smooth_at_a_point_for_a_finite_type_scheme_over_a_field}{smooth}, is \CrefAndHyperrefIfExist{definition:projective_variety_over_a_field}{projective}, and has \CrefAndHyperrefIfExist{definition:geometrically_connected_scheme_over_a_field}{geometrically connected} \CrefAndHyperrefIfExist{definition:connected_component_of_a_point_of_a_topological_space}{connected components}.
        The Jacobian is an \CrefAndHyperrefIfExist{definition:abelian_variety_over_a_field}{abelian variety} over $k$ and may be called the \hldef{Jacobian variety of $X$}. It may be denoted by notations such as \hl{$\Jac(X)$}, \hl{$\Jac_X$}, or \hl{$J_X$}. 
        
        Under these assumptions, note that $\Jac(X)$ coincides with the \CrefAndHyperrefIfExist{definition:picard_variety_of_a_smooth_projective_geometrically_integral_variety_over_a_field}{Picard variety} of $X/k$. Equivalently, $J_X$ parametrizes isomorphism classes of degree-zero \CrefAndHyperrefIfExist{definition:free_and_locally_free_modules_on_a_ringed_site}{line bundles} on $X$, where a line bundle $\mathcal{L}$ has \hldef{degree zero} if its degree (computed as the sum of orders of any rational function defining it) vanishes. 

        % Let $X$ be a \CrefAndHyperrefIfExist{definition:smooth_at_a_point_for_a_finite_type_scheme_over_a_field}{smooth}, \CrefAndHyperrefIfExist{definition:projective_variety_over_a_field}{projective}, \CrefAndHyperrefIfExist{definition:disconnected_and_connected_topological_spaces}{connected} \CrefAndHyperrefIfExist{definition:algebraic_curve_over_a_field}{curve} of \CrefAndHyperrefIfExist{definition:geometric_genus_of_a_smooth_projective_variety_over_a_field}{genus} $g \geq 1$ over a \CrefAndHyperrefIfExist{definition:field}{field} $k$.


        The Jacobian variety $\Jac(X)$ has \CrefAndHyperrefIfExist{definition:dimension_of_a_scheme}{dimension} $g$ where $g$ is the \CrefAndHyperrefIfExist{definition:geometric_genus_of_a_smooth_projective_variety_over_a_field}{geometric genus} of $X$.

        When $X$ is equipped with a base point, the Jacobian variety carries a canonical \CrefAndHyperrefIfExist{definition:principal_polarization_of_an_abelian_variety_over_a_field}{principal polarization} $\Theta$, making $(J_X, \Theta)$ a principally polarized abelian variety.
    \end{enumerate}

\end{definition}

\section{N\'eron models of abelian varieties}


\begin{definition}[Dedekind domain] \label{definition:dedekind_domain}
An \CrefAndHyperrefIfExist{definition:domain_integral_domain}{integral domain} $R$ is called a \hldef{Dedekind domain} or a \hldef{Dedekind ring} if it satisfies the following equivalent conditions:
\begin{itemize}
    \item $R$ is \CrefAndHyperrefIfExist{definition:noetherian_ring}{Noetherian}, \CrefAndHyperrefIfExist{definition:integral_element_over_a_ring}{integrally closed} in its \CrefAndHyperrefIfExist{definition:field_of_fractions_of_an_integral_domain}{field of fractions}, and every nonzero \CrefAndHyperrefIfExist{definition:prime_and_maximal_ideal_of_a_ring}{prime ideal} of $R$ is \CrefAndHyperrefIfExist{definition:prime_and_maximal_ideal_of_a_ring}{maximal}.

    \item Equivalently: for every nonzero prime ideal $\mathfrak{p}$ of $R$, the \CrefAndHyperrefIfExist{definition:localization_of_a_commutative_ring_by_a_multiplicative_subset}{localization} $R_{\mathfrak{p}}$ is a \CrefAndHyperrefIfExist{definition:discrete_valuation_ring}{discrete valuation ring}.
\end{itemize}
\end{definition}

\begin{definition}[Dedekind scheme] \label{definition:dedekind_scheme}
    A \hldef{Dedekind scheme} is a \CrefAndHyperrefIfExist{definition:scheme}{scheme} $S$ that is \CrefAndHyperrefIfExist{definition:integral_scheme}{integral}, \CrefAndHyperrefIfExist{definition:locally_noetherian_and_noetherian_scheme}{Noetherian}, \CrefAndHyperrefIfExist{definition:normal_and_factorial_scheme}{normal}, and of \CrefAndHyperrefIfExist{definition:dimension_of_a_topological_space}{Krull dimension} $1$. Equivalently, every point $s \in S$ has an open neighborhood $U = \Spec R$ such that $R$ is a \CrefAndHyperrefIfExist{definition:dedekind_domain}{Dedekind domain}. In particular, for every point $s \in S$, the local ring $\mathcal{O}_{S, s}$ is either a field (if $s$ is the generic point) or a \CrefAndHyperrefIfExist{definition:discrete_valuation_ring}{discrete valuation ring} (if $s$ is a closed point).
\end{definition}

\begin{definition}[Special fiber and generic fiber] \label{definition:special_and_generic_fiber_of_a_scheme_over_a_dedekind_scheme}
    Let $S$ be a \CrefAndHyperrefIfExist{definition:integral_scheme}{integral} \CrefAndHyperrefIfExist{definition:scheme}{scheme}, let $\eta$ be its \CrefAndHyperrefIfExist{definition:generic_point_of_a_topological_space}{generic point}, and let $K = \mathcal{O}_{S, \eta}$ be the \CrefAndHyperrefIfExist{definition:function_field_of_an_integral_scheme}{function field} of $S$. Let $s \in S$ be a closed point.

    \begin{enumerate}
        \item For any \CrefAndHyperrefIfExist{definition:morphism_of_schemes}{morphism of schemes} $f: X \to S$, the \hldef{generic fiber of $X$ over $S$} (or simply the \hldef{generic fiber}) is the \CrefAndHyperrefIfExist{definition:base_change_of_a_morphism_in_a_category_by_a_morphism}{base change}:
        $$\hlin{X_\eta = X \times_S \Spec K}$$
        of $f$ along $\Spec K \to S$.

        \item For any morphism of schemes $f: X \to S$, the \hldef{special fiber of $X$ at $s$} is the \CrefAndHyperrefIfExist{definition:base_change_of_a_morphism_in_a_category_by_a_morphism}{base change}:
        $$\hlin{X_s = X \times_S \Spec k(s)}$$
        of $f$ along $\Spec k(s) \to S$, where $k(s) = \mathcal{O}_{S, s} / \mathfrak{m}_s$ is the \CrefAndHyperrefIfExist{definition:residue_field_of_a_scheme_at_a_point}{residue field} at $s$.

        In the case that $S$ has only closed point (e.g. $S$ is the \CrefAndHyperrefIfExist{definition:ring_of_integers_of_a_global_or_local_ring}{ring of integers} of a \CrefAndHyperrefIfExist{definition:local_field}{local field}), the \hldef{special fiber of $X$} refers to the special fiber of $X$ at this unique closed point.
    \end{enumerate}
\end{definition}

\begin{definition}[N\'eron model of an abelian variety] \label{definition:neron_model_of_an_abelian_variety_over_the_function_field_of_an_integral_scheme.tex}
    Let $S$ be an \CrefAndHyperrefIfExist{definition:integral_scheme}{integral scheme}, let $\eta$ be its \CrefAndHyperrefIfExist{definition:generic_point_of_a_topological_space}{generic point}, and let $K = \mathcal{O}_{S, \eta}$ be the \CrefAndHyperrefIfExist{definition:function_field_over_a_field}{function field} of $S$. Let $A_K$ be an \CrefAndHyperref{definition:abelian_variety_over_a_field}{abelian variety over $K$}.
    % Let $S$ be a \CrefAndHyperrefIfExist{definition:dedekind_scheme}{Dedekind scheme}, let $\eta$ be its \CrefAndHyperrefIfExist{definition:generic_point_of_a_topological_space}{generic point}, and let $K = \mathcal{O}_{S, \eta}$ be the \CrefAndHyperrefIfExist{definition:function_field_over_a_field}{function field} of $S$. Let $A_K$ be an \CrefAndHyperref{definition:abelian_variety_over_a_field}{abelian variety over $K$}.

    A \hldef{N\'eron model of $A_K$ over $S$} is a \CrefAndHyperrefIfExist{definition:smooth_morphism_of_schemes}{smooth} \CrefAndHyperrefIfExist{definition:quasi_separated_morphism_of_schemes}{separated} \CrefAndHyperrefIfExist{definition:group_scheme_over_a_scheme}{group scheme} $\mathcal{A}$ of \CrefAndHyperrefIfExist{definition:finite_type_morphism_of_schemes}{finite type} over $S$ whose \CrefAndHyperrefIfExist{definition:special_and_generic_fiber_of_a_scheme_over_a_dedekind_scheme}{generic fiber} $\mathcal{A}_\eta$ is isomorphic to $A_K$, satisfying the following universal property (the \hldef{N\'eron mapping property}): for every smooth $S$-scheme $U$ and every $K$-morphism $f_K: U_\eta \to A_K$, there exists a unique $S$-morphism $f: U \to \mathcal{A}$ extending $f_K$.

    If such a model exists, it is unique up to unique isomorphism.
\end{definition}


\section{Galois cohomology}
\begin{definition} \label{definition:group_cohomology_homology_of_a_module_over_a_group_ring_over_a_ring}
Let $G$ be a \CrefAndHyperrefIfExist{definition:group}{group} and $R$ a \CrefAndHyperrefIfExist{definition:ring}{not necessarily commutative ring}. Let $A$ be a \CrefAndHyperrefIfExist{definition:left_right_module_over_a_group_ring_of_a_group_over_a_ring}{left/right $R[G]$-module}.

\begin{enumerate}
    \item The \hldef{cohomology group of $G$ with coefficients in $A$} or the \hldef{group cohomology of $A$ as an $R[G]$-module}, denoted \hl{$H^n(G, A)$} for $n \geq 0$, is the $n$-th \CrefAndHyperrefIfExist{definition:left_right_derived_functors_of_a_right_left_exact_functor_between_abelian_categories_where_source_has_enough_projectives_injectives}{right derived functor} object $(R^n(-^G))(A)$ of the \CrefAndHyperrefIfExist{definition:invariant_subgroup_of_a_group_acting_on_an_abelian_group}{invariants} functor $A \mapsto A^G$, which is \CrefAndHyperrefIfExist{definition:exact_functor_between_abelian_categories}{left exact} by \Cref{lemma:invariant_coinvariant_functors_for_group_rings_are_adjoint_to_trivial_module_functor}.

    \item The \hldef{homology group of $G$ with coefficients in $A$} or the \hldef{group homology of $A$ as an $R[G]$-module}, denoted \hl{$H_n(G, A)$} for $n \geq 0$, is the $n$-th \CrefAndHyperrefIfExist{definition:left_right_derived_functors_of_a_right_left_exact_functor_between_abelian_categories_where_source_has_enough_projectives_injectives}{left derived functor} object $(L_n(-_G))(A)$ of the \CrefAndHyperrefIfExist{definition:invariant_subgroup_of_a_group_acting_on_an_abelian_group}{coinvariants} functor $A \mapsto A_G$, which is \CrefAndHyperrefIfExist{definition:exact_functor_between_abelian_categories}{right exact} by \Cref{lemma:invariant_coinvariant_functors_for_group_rings_are_adjoint_to_trivial_module_functor}.
\end{enumerate}
 
\end{definition}

\begin{definition} \label{definition:continuous_cochain_complex_of_a_topological_module_over_a_group_ring_of_a_topological_group_over_a_topological_ring}
Let $R$ be a \CrefAndHyperrefIfExist{definition:topological_ring}{topological ring} and $G$ a \CrefAndHyperrefIfExist{definition:topological_group}{topological group}. Let $R\text{-}\mathbf{TopMod}_G$ be the \CrefAndHyperrefIfExist{definition:category_of_topological_modules_over_a_group_ring_of_a_topological_group_over_a_topological_ring}{category of (left) topological $G$-modules over $R$}.

Let $M$ be an object in $R\text{-}\mathbf{TopMod}_G$. The \hldef{continuous cochain complex} \hl{$C^\bullet_{\text{cont}}(G, M)$} is defined as follows:
\begin{enumerate}
    \item $C^n_{\text{cont}}(G, M) = \text{Map}_{\text{cont}}(G^n, M)$, the $R$-module of continuous maps from $G^n$ to $M$.
    \item The \hldef{coboundary map} \hl{$d^n: C^n_{\text{cont}}(G, M) \to C^{n+1}_{\text{cont}}(G, M)$} is given by:
    $$(d^n f)(g_1, \dots, g_{n+1}) = g_1 f(g_2, \dots, g_{n+1}) + \sum_{i=1}^n (-1)^i f(g_1, \dots, g_i g_{i+1}, \dots, g_{n+1}) + (-1)^{n+1} f(g_1, \dots, g_n)$$
\end{enumerate}
\end{definition}
\begin{definition} \label{definition:continuous_cohomology_groups_of_a_topological_module_over_a_topological_group_and_a_topological_ring}

    % \TODO{It is probably more accurate to a priori define continuous cohomology by the chain complexes and then define it as the derived functor assuming that enough injectives existe}


Let $R$ be a \CrefAndHyperrefIfExist{definition:topological_ring}{topological ring} and $G$ a \CrefAndHyperrefIfExist{definition:topological_group}{topological group}. Let $R\text{-}\mathbf{TopMod}_G$ be the \CrefAndHyperrefIfExist{definition:category_of_topological_modules_over_a_group_ring_of_a_topological_group_over_a_topological_ring}{category of (left) topological $G$-modules over $R$}.

The \hldef{continuous cohomology groups} \hl{$H^n_{\text{cont}}(G, M)$}  are defined as the cohomology groups of the \CrefAndHyperref{definition:continuous_cochain_complex_of_a_topological_module_over_a_group_ring_of_a_topological_group_over_a_topological_ring}{continuous cochain complex} $C^\bullet_{\text{cont}}(G, M)$.

The continuous cohomology group is also denoted by \hl{$H^n(G, M)$} if the continuity is clear from context, but this should not be confused for the \CrefAndHyperrefIfExist{definition:group_cohomology_homology_of_a_module_over_a_group_ring_over_a_ring}{(non-continuous) cohomology group} $H^n(G,M)$ of $M$ as a $G$-module (without consideration of the topological structures).

\end{definition}
\begin{theorem} \label{theorem:continuous_cohomology_of_topological_module_over_topological_gropu_and_topological_ring_can_be_computed_by_continuous_cohomology_groups}
    \TODO{verify these facts}

Let $R$ be a \CrefAndHyperrefIfExist{definition:topological_ring}{topological ring} and $G$ a \CrefAndHyperrefIfExist{definition:topological_group}{topological group}. Let $R\text{-}\mathbf{TopMod}_G$ be the \CrefAndHyperrefIfExist{definition:category_of_topological_modules_over_a_group_ring_of_a_topological_group_over_a_topological_ring}{category of (left) topological $G$-modules over $R$}.

Consider the \CrefAndHyperrefIfExist{definition:G_invariant_functor_of_topological_group_module_over_topological_ring}{$G$-invariant functor} $(-)^G: R\text{-}\mathbf{TopMod}_G \to R\text{-}\mathbf{TopMod}$ (\Cref{definition:categories_of_topological_modules_over_topological_rings}).

Assume either of the following:
\begin{enumerate}
    \item $G$ is \CrefAndHyperrefIfExist{definition:profinite_group}{profinite}, $R$ is discrete, and $M$ is \CrefAndHyperrefIfExist{definition:discrete_group}{discrete}; in this case, by the ``right derived functors of $(-)^G$'', we mean the right derived functor on the category of discrete $G$-modules, which has enough injectives \TODO{cite; apparently the category of discrete $G$-modules over a topological ring $R$ where $G$ is profinite is an abelian category with enough injectives}.
    \item $G$ is locally compact, and $M$ is complete; in this case, by the ``right derived functors of $(-)^G$'', we mean the right derived functor in the sense of relative homological algebra, using resolutions that split as sequences of topological spaces \TODO{figureo ut what this means}.
\end{enumerate}

The \CrefAndHyperrefIfExist{definition:continuous_cohomology_groups_of_a_topological_module_over_a_topological_group_and_a_topological_ring}{continuous cohomology groups} $H^n_{\text{cont}}(G,M)$, i.e. the  \CrefAndHyperrefIfExist{definition:homology_and_cohomology_objects_for_a_chain_complex_in_an_additive_category}{cohomology} of the \CrefAndHyperrefIfExist{definition:continuous_cochain_complex_of_a_topological_module_over_a_group_ring_of_a_topological_group_over_a_topological_ring}{continuous cochain complex} $C^\bullet_{\text{cont}}(G, M)$, is naturally isomorphic to the the right derived functors of $(-)^G$. In other words,
$$H^n(C^\bullet_{\text{cont}}(G, M)) \cong (R^n(-^G))(M)$$
\end{theorem}
\begin{definition} \label{definition:galois_cohomology_of_a_galois_field_extension_with_coefficients_in_a_topological_galois_module_over_a_topological_ring}
    Let $R$ be a \CrefAndHyperrefIfExist{definition:topological_ring}{topological ring} (e.g., a $\ell$-adic \CrefAndHyperrefIfExist{definition:ring}{ring} $\mathbb{Z}_\ell$) \TODO{ell-adic ring}.
    \begin{enumerate}
        \item Let $L/K$ be a \CrefAndHyperrefIfExist{definition:galois_extension_of_fields}{Galois field extension}.   Let $M$ be a \CrefAndHyperrefIfExist{definition:topological_module_over_a_group_ring_of_a_topological_gropu_over_a_topological_ring}{topological $G$-module} over $R$

        The \hldef{Galois cohomology group of $L/K$ with coefficients in $M$} is:
        $$\hlin{H^n(L/K, M) = H^n_{\text{cont}}(\Gal(L/K), M)}$$
        (\Cref{definition:continuous_cohomology_groups_of_a_topological_module_over_a_topological_group_and_a_topological_ring}).

        \TODO{absolute galois group}
        \item Let $K$ be a \CrefAndHyperrefIfExist{definition:field}{field} with absolute Galois group $G_K = \Gal(\overline{K}/K)$. Let $M$ be a \CrefAndHyperrefIfExist{definition:topological_module_over_a_group_ring_of_a_topological_gropu_over_a_topological_ring}{topological $G$-module} over $R$.

        The \hldef{Galois cohomology group of $K$ with coefficients in $M$} is:
        $$\hlin{H^n(K, M) = H^n(\overline{K}/K, M) = H^n_{\text{cont}}(\Gal(\overline{K}/K), M)}.$$
    \end{enumerate}

This framework is used to define $\ell$-adic cohomology, where $M$ is often a $\mathbb{Z}_\ell$-module or a $\mathbb{Q}_\ell$-vector space arising from the Tate module of an abelian variety or the étale cohomology of a variety.
\end{definition}

\section{Selmer groups}
\begin{definition}[Kummer map associated to an isogeny of an abelian variety] \label{definition:kummer_map_associated_to_an_isogeny_of_a_commutative_group_schemes_over_a_field}
    Let $k$ be a \CrefAndHyperrefIfExist{definition:field}{field}, and let $A$ and $B$ be \CrefAndHyperrefIfExist{definition:algebraic_group_scheme_over_a_scheme}{commutative group schemes} over $S$.  
    Suppose $\varphi: A \to B$ is an \CrefAndHyperrefIfExist{definition:isogeny_of_algebraic_group_schemes_over_a_scheme}{isogeny over $S$}.  
    Denote by $G_k = \mathrm{Gal}(\overline{k}/k)$ the \CrefAndHyperrefIfExist{definition:absolute_galois_group_of_a_field}{absolute Galois group} of $k$.  
    % Let $k$ be a \CrefAndHyperrefIfExist{definition:field}{field}, and let $A$ and $B$ be \CrefAndHyperrefIfExist{definition:abelian_variety_over_a_field}{abelian varieties} defined over $k$.  
    % Suppose $\varphi: A \to B$ is an \CrefAndHyperrefIfExist{definition:isogeny_of_abelian_varieties_over_a_field}{isogeny defined over $k$}.  
    % Denote by $G_k = \mathrm{Gal}(\overline{k}/k)$ the \CrefAndHyperrefIfExist{definition:absolute_galois_group_of_a_field}{absolute Galois group} of $k$.  

    The \CrefAndHyperrefIfExist{definition:exact_sequence_short_exact_sequence_of_group_homomorphisms}{short exact sequence} 
    $$0 \to \ker \varphi \to A(\overline{k}) \xrightarrow{\varphi} B(\overline{k}) \to 0.$$
    yields the following \CrefAndHyperrefIfExist{theorem:long_exact_sequence_in_homology_cohomology_of_a_short_exact_sequence_of_chain_complexes_in_an_abelian_category}{long exact sequence} in \CrefAndHyperrefIfExist{definition:galois_cohomology_of_a_galois_field_extension_with_coefficients_in_a_topological_galois_module_over_a_topological_ring}{Galois cohomology}:
    $$
    \begin{aligned}
    0 \to \ker \varphi(k) \to &A(k) \xrightarrow{\varphi} B(k) \\
    \xrightarrow{\delta_\varphi} H^1(G_k, \ker \varphi) \to &H^1(G_k, A(\overline{k})) \to H^1(G_k, B(\overline{k})) \to \cdots.
    \end{aligned}
    $$

    The \hldef{Kummer map associated to the isogeny $\varphi$} is the Galois cohomological connecting homomorphism
    $$\delta_\varphi: B(k) \to H^1(G_k, \ker \varphi),$$
    above. In fact, there is a short exact sequence
    $$0 \to B(k) / \varphi(A(k)) \xrightarrow{\delta'_{\varphi}} H^1(G_{k}, \ker \varphi) \to H^1(G_{k}, A)[\varphi] \to 0$$
    where the map $\delta'_{\varphi}$ is induced by $\delta_{\varphi}$. We may also let the \hldef{Kummer map associated to $\varphi$} refer to this map $\delta'_{\varphi}$. We may also use the abuse of notation \hl{$\delta_\varphi$} to denote $\delta'_{\varphi}$.
    % The map $\delta_\varphi$ measures the obstruction to lifting points of $B(k)$ to points of $A(k)$ through the isogeny $\varphi$, fundamental in descent theory.

\end{definition}

\begin{definition}[Selmer group for an isogeny of abelian group schemes over a global field] \label{definition:selmer_group_for_an_isogeny_of_commutative_group_schemes_over_a_global_field}
    \TODO{It should be possible to define this for more general group schemes}
    Let $K$ be a \CrefAndHyperrefIfExist{definition:global_field}{global field}, and let $A$ and $B$ be \CrefAndHyperrefIfExist{definition:algebraic_group_scheme_over_a_scheme}{abelian group schemes} defined over $K$.  
    Suppose $\varphi: A \to B$ is an \CrefAndHyperrefIfExist{definition:isogeny_of_algebraic_group_schemes_over_a_scheme}{isogeny defined over $K$}.  

    Denote by $G_K = \mathrm{Gal}(\overline{K}/K)$ the \CrefAndHyperrefIfExist{definition:absolute_galois_group_of_a_field}{absolute Galois group} of $K$. For each \CrefAndHyperrefIfExist{definition:place_of_a_global_field}{place} $v$ of $K$, fix an extension $v$ to $\barK$, which yields an embedding $\barK \subset \barK_v$ and a \CrefAndHyperrefIfExist{definition:decomposition_group_of_a_general_galois_extension_of_valuation_fields}{decomposition group} $G_v \subset G_K$.
    
    The \hldef{Selmer group of $\varphi$ over $K$}, denoted by notations such as \hl{$\mathrm{Sel}(A/K)$}, \hl{$\mathrm{Sel}^\varphi(A/K)$}, \hl{$\mathrm{Sel}^{(\varphi)}(A/K)$}, \hl{$\mathrm{Sel}_\varphi(A/K)$}, is the subgroup of the \CrefAndHyperrefIfExist{definition:galois_cohomology_of_a_galois_field_extension_with_coefficients_in_a_topological_galois_module_over_a_topological_ring}{Galois cohomology group} $H^1(K, \ker \varphi)$ given by
    $$\mathrm{Sel}^\varphi(A/K) := \ker \left( H^1(G_K, A[\varphi]) \to \prod_v H^1(G_v, A(\barK))[\varphi] \right),$$
    \TODO{define the kummer map associated to $\varphi$}
    where the product runs over all \CrefAndHyperrefIfExist{definition:place_of_a_global_field}{places} $v$ of $K$, and $H^1(G_v, A)[\varphi]$ denotes the image of the local Kummer map associated to $\varphi$.  

    We describe the map 
    \begin{equation} \label{eq:map_whose_kernel_is_selmer_group}
    H^1(G_K, A[\varphi]) \to \prod_v H^1(G_v, A(\barK))[\varphi]
    \end{equation}
    \TODO{define a decomposition group}
    used to define the kernel above: 
     Note that $G_v$ acts on $A(\barK_v)$ and $B(\barK_v)$, and the base change $\varphi_v$ of $\varphi$ to $K_v$ induces a \CrefAndHyperrefIfExist{definition:kummer_map_associated_to_an_isogeny_of_a_commutative_group_schemes_over_a_field}{Kummer short exact sequence}
    $$0 \to B(K_v) / \varphi(A(K_v)) \xrightarrow{\delta_{\varphi_v}} H^1(G_v, A[\varphi]) \to H^1(G_v, A(\barK_v))[\varphi] \to 0.$$
    The natural inclusions $G_v \subset G_{K}$ and $E(K) \subset E(K_v)$ give \CrefAndHyperrefIfExist{definition:restriction_map_of_continuous_group_cohomology_of_a_topological_group_on_a_topological_module}{restriction maps} $H^1(G_K, A[\varphi]) \to H^1(G_v, A[\varphi])$ on Galois cohomology, and we have a commutative diagram
    \begin{center}
        \begin{tikzcd}
        0 \ar[r] & B(K) / \varphi(A(K))  \ar[r, "\delta_{\varphi}"] \ar[d] & H^1(G_K, A[\varphi])\ar[r] \ar[d, "\prod_v \operatorname{res}_v" ] \ar[rd, dotted] & H^1(G_K,A(\barK))[\phi] \ar[r] \ar[d] & 0 \\
        0 \ar[r] & \prod_v B(K_v) / \varphi(A(K_v)) \ar[r, "\delta_{\varphi_v}"]  & \prod_v H^1(G_v,A[\varphi])  \ar[r] & \prod_v H^1(G_v,A(\barK_v))[\phi] \ar[r] & 0
        \end{tikzcd}
    \end{center}
    The map \eqref{eq:map_whose_kernel_is_selmer_group} is the one given in the above commutative diagram by the dotted arrow. 

    Equivalently, 
    $$\Sel^\varphi(A/K) = \left\{c \in H^1(K, A[\varphi]): \text{ for all places } v \text{ of } K, \operatorname{res}_v(c) \in \operatorname{Im} \delta_{\varphi_v} \right \}$$
    \CrefIfExists{proposition:selmer_group_of_a_commutative_group_scheme_over_a_global_field_is_group_of_galois_cocycles_whose_restriction_are_local_mordell_weil_groups}
\end{definition}


\section{\'Etale sheaves}

\input{../_concepts/theorem_locally_constant_sheaves_on_X_etale_are_equivalent_to_representations_of_the_etale_fundamental_group.tex}
\begin{theorem} \label{theorem:lisse_ell_adic_sheaves_on_connected_noetherian_schemes_are_equivalent_to_representations_of_etale_fundamental_group_over_adic_ring}
    Let $(R, \lambda)$ be a \CrefAndHyperrefIfExist{definition:complete_ring_with_respect_to_an_ideal}{complete} \CrefAndHyperrefIfExist{definition:discrete_valuation_ring}{DVR} such that $R/\lambda^n$ is finite for all $n \geq 1$. Let $E$ be the \CrefAndHyperrefIfExist{definition:field_of_fractions_of_an_integral_domain}{fraction field} of $R$. Let $K$ be an algebraic extension of $E$. Let $X$ be a \CrefAndHyperrefIfExist{definition:disconnected_and_connected_topological_spaces}{connected} \CrefAndHyperrefIfExist{definition:locally_noetherian_and_noetherian_scheme}{noetherian scheme}. Let $\barx \in X$ be a \CrefAndHyperrefIfExist{definition:geometric_point_of_a_scheme}{geometric point}.

    \begin{enumerate}
        \item There is an equivalence of categories between:
        \[
        \left\{
        \begin{matrix}
            \text{lisse } \lambda\text{-adic sheaves on } X_{\acute{e}tale}
        \end{matrix}
        \right\}
        \longleftrightarrow
        \left\{
        \begin{matrix}
            \text{finitely generated free } R\text{-modules } M \\
            \text{with a continuous } \pioneet(X, \overline{x})\text{-action}
        \end{matrix}
        \right\}.
        \]
        \CrefIfExists{definition:lambda_adic_sheaf_on_the_small_etale_site_of_a_scheme_for_an_ideal_of_a_commutative_ring}
        \CrefIfExists{definition:lisse_sheaf_on_a_scheme}
        \CrefIfExists{definition:continuous_representation_of_a_topological_group_on_a_topological_module_over_a_topological_ring}

        \item There is an equivalence of categories between:
        \[
        \left\{
        \begin{matrix}
            \text{lisse objects of } \Sh(X, E)
        \end{matrix}
        \right\}
        \longleftrightarrow
        \left\{
        \begin{matrix}
            \text{finite-dimensional } E\text{-vector spaces } V \\
            \text{with a continuous } \pioneet(X, \overline{x})\text{-action}
        \end{matrix}
        \right\}.
        \]
        \CrefIfExists{definition:category_of_etale_sheaves_with_rational_coefficients_in_the_fraction_field_of_a_complete_dvr}
        \CrefIfExists{definition:lisse_sheaf_on_a_scheme}

        \item Let $K$ be an algebraic extension of $E$. There is an equivalence of categories between:
        \[
        \left\{
        \begin{matrix}
            \text{lisse objects of } \Sh(X, K)
        \end{matrix}
        \right\}
        \longleftrightarrow
        \left\{
        \begin{matrix}
            \text{finite-dimensional } K\text{-vector spaces } V \\
            \text{with a continuous } \pioneet(X, \overline{x})\text{-action}
        \end{matrix}
        \right\}.
        \]
        \CrefIfExists{definition:category_of_etale_sheaves_with_rational_coefficients_in_an_algebraic_extension_of_the_fraction_field_of_a_complete_dvr}
        \CrefIfExists{definition:lisse_sheaf_on_a_scheme}
    \end{enumerate}

\end{theorem}

\begin{definition} \label{definition:decomposition_group_of_a_general_galois_extension_of_valuation_fields}
    Assume \Cref{context:extension_of_valuation_fields} and further assume that $L/K$ is a \CrefAndHyperrefIfExist{definition:galois_extension_of_fields}{Galois extension}.

    The \hldef{decomposition group of $L/K$} is the subgroup 
    $$\hlin{D_w = \{\sigma \in \Gal(L/K) : w \circ \sigma = w\}}.$$
    \CrefIfExists{definition:galois_extension_of_fields} It is also the \CrefAndHyperrefIfExist{definition:projective_and_inductive_limits_in_categories}{inverse limit} of the \CrefAndHyperrefIfExist{definition:decomposition_group_of_a_finite_galois_extension_of_valuation_fields}{decomposition groups} of the finite Galois subextensions $E/K$ within $L$ (with valuation restricted from the valuation of $L$). It is a closed subgroup of $\Gal(L/K)$.

    Letting $(K,v)$ be a valuation field, its \hldef{decomposition group} \hl{$D_{K,v}$} (also sometimes denoted by \hl{$G_{K,v}$}) is the decomposition group of $K^{\text{sep}} / K$ in the above sense where a valuation on $K^{\text{sep}}$ extending $v$ is fixed.

\end{definition}
\begin{definition} \label{definition:inertia_group_of_a_general_galois_extension_of_valuation_fields}
    Assume \Cref{context:extension_of_valuation_fields} and further assume that $L/K$ is a \CrefAndHyperrefIfExist{definition:galois_extension_of_fields}{Galois extension}.

    There is a natural surjective homomorphism 
    $$D_w \to \Gal(k_w / k_v)$$
    (\Cref{definition:decomposition_group_of_a_general_galois_extension_of_valuation_fields})
    on the \CrefAndHyperrefIfExist{definition:residue_field_of_a_ring_at_a_point}{residue fields}. The \hldef{inertia group of $L/K$} is its kernel: 
    $$\hlin{I_w = \{\sigma \in D_w : \sigma(x) \equiv x \pmod{\mathfrak{m}_v} \text{ for all } x \in \mathcal{O}_v \}}.$$
    (\Cref{definition:decomposition_group_of_a_general_galois_extension_of_valuation_fields}) \CrefIfExists{definition:galois_extension_of_fields} It is also the \CrefAndHyperrefIfExist{definition:projective_and_inductive_limits_in_categories}{inverse limit} of the \CrefAndHyperrefIfExist{definition:decomposition_group_of_a_finite_galois_extension_of_valuation_fields}{decomposition groups} of the finite Galois subextensions $E/K$ within $L$ (with valuation restricted from the valuation of $L$). It is a closed subgroup of $D_w$.

    Letting $(K,v)$ be a valuation field, its \hldef{inertia group} \hl{$I_{K,v}$} is the inertia group of $K^{\text{sep}} / K$ in the above sense where a valuation on $K^{\text{sep}}$ extending $v$ is fixed.

\end{definition}

\begin{proposition}[see {\cite[Tag 0BQM]{stacks-project}}] \label{proposition:mor_from_et_fund_grp_of_gen_pt_of_norm_int_sch_to_et_fund_grp_of_scheme_is_quo_map_of_gal_grp_of_sep_clos_of_func_field_to_gal_grp_of_max_unr_ext}
    Let $X$ be a \CrefAndHyperrefIfExist{definition:normal_and_factorial_scheme}{normal} \CrefAndHyperrefIfExist{definition:integral_scheme}{integral scheme} \CrefAndHyperrefIfExist{definition:function_field_over_a_field}{function field} $K$. Let $\eta$ be the generic point of $X$ and let $\overline{\eta}$ be a \CrefAndHyperrefIfExist{definition:geometric_point_of_a_scheme}{geometric point} of $X$ whose image is $\eta$. 
    The canonical continuous group homomorphism
    $$\Gal(K^\text{sep} / K) \cong \pioneet(\eta, \overline{\eta}) \to \pi_1(X, \overline{\eta})$$
    \CrefIfExists{proposition:galois_group_of_a_galois_extension_of_fields_is_profinite}\CrefIfExists{definition:separable_closure_of_a_field}\CrefIfExists{definition:etale_fundamental_group_of_a_connected_scheme}
    \TODO{better describe the canonical map}

    is identified with the quotient map 
    $$\Gal(K^\text{sep} / K) \to \Gal(M/K)$$
    \TODO{defined unramifiedness of $X$ in $L$}
    where $M \subseteq K^{\text{sep}}$ is the union of the finite subextensions $L$ such that $X$ is unramified in $L$. 
\end{proposition}
\begin{definition} \label{definition:drop_of_a_lisse_sheaf_on_dense_open_of_integral_normal_scheme_at_point_in_complement}
\TODO{AI generated, verify}
Let $X$ be an \CrefAndHyperrefIfExist{definition:integral_scheme}{integral} \CrefAndHyperrefIfExist{definition:normal_and_factorial_scheme}{normal scheme} with \CrefAndHyperrefIfExist{definition:generic_point_of_a_topological_space}{generic point} $\eta$. Let $\overline{\eta}$ be a geometric point of $X$ whose image is $\eta$. Let $U \subseteq X$ be a non-empty \CrefAndHyperrefIfExist{definition:open_and_closed_subscheme_of_a_scheme}{open subscheme}. Let $x \in X \setminus U$ be a point. Let $I_x$ be the \CrefAndHyperrefIfExist{definition:decomposition_and_inertia_subgroups_of_etale_fundamental_group_of_dense_open_of_integral_normal_scheme_at_point_in_complement}{inertia subgroup} of $\pioneet(U, \overline{\eta})$ associated to $x$, well defined up to conjugation.

Let $\Lambda$ be a finite ring and let $\calF$ be a locally constant sheaf of finitely generated $\Lambda$-modules on $U_{\et}$. The inertia subgroup $I_x$ acts continuously on the stalk $\calF_{\overline{\eta}}$, which is a finitely generated $\Lambda$-module. For any field $k$, let $V = \calF_{\overline{\eta}} \otimes_\Lambda k$ be the resulting finite-dimensional $k$-vector space with induced $I_x$-action.

The \hldef{drop of $\calF$ at $x$} is the \CrefAndHyperrefIfExist{definition:drop_of_a_representation_of_a_group_on_a_vector_space}{drop of the representation} $V$ under the action of $I_x$:
$$\hlin{\operatorname{drop}_x(\calF) = \dim_k(V) - \dim_k(V^{I_x})}.$$
When $\Lambda$ is a field, this simplifies to $\dim_\Lambda(\calF_{\overline{\eta}}) - \dim_\Lambda(\calF_{\overline{\eta}}^{I_x})$.

For a \CrefAndHyperrefIfExist{definition:lisse_sheaf_on_a_scheme}{lisse $R$-sheaf} $\calF = (\calF_n, u_n)_{n \geq 1}$ on $U_{\et}$ where $(R, \lambda)$ is a complete DVR with fraction field $E$, the \hldef{drop of $\calF$ at $x$} is defined as the drop of the $E$-vector space $\calF_{\overline{\eta}} \otimes_R E$ under the induced $I_x$-action.

The drop is independent of the choice of valuation centered at $x$ because different choices yield conjugate inertia subgroups, and conjugation does not change the dimension of the invariant subspace. The drop vanishes if and only if $\calF$ is \CrefAndHyperrefIfExist{definition:ramification_of_an_etale_sheaf_on_a_connected_normal_scheme_at_a_geometric_point}{unramified at $x$}.
\end{definition}
